\documentclass[twoside]{article}
\usepackage[paperheight=395pt, paperwidth=240pt, width=166pt, height=308pt, top=42pt, headsep=2pt]{geometry}
\usepackage[noinfo, frame, a6, center, lualatex]{crop}
\usepackage{shandean}
\usepackage{graphicx}
%
\title{A Sentimental Journey through France and Italy}
\author{Laurence Sterne}
\date{1768}
\begin{document}
\def\vol{II}
\pagenumbering{roman}
\pagestyle{empty}
\centerline{A}
\centerline{\tstrut 28pt\large \lsv{24}{SENTIMENTAL\enspace JOURNEY}}
\centerline{\tstrut 24pt\small \lsv{36}{\textsc{through}}}
\centerline{\tstrut 28pt\lss{\Large FRANCE}\quad\lss{AND}\quad\lss{\Large ITALY}.}
\centerline{\tstrut 32pt\small \lsv{36}{\textsc{by}}}
\centerline{\tstrut 22pt\large \lsv{36}{\textsc{Mr}}.\quad \lsv{48}{YORICK}.}

\vskip 36pt
\hrule
\centerline{\tstrut 16pt\large\lsv{44}{VOL}.\quad II.}
\vskip 8pt
\hrule
\vskip 2pt
\hrule
\centerline{\tstrut 54pt\small \lsv{72}{LONDON}:}
\centerline{\tstrut 18pt\smaller Printed for T. \lsv{18}{\scshape Becket}\ and\ P.\@
A.\@ \lsv{18}{\scshape De Hondt},}
\centerline{\tstrut 10pt\smaller in the Strand.\enspace MDCCLXVIII.}
\newpage
\null
\newpage
\null\vskip 1in
\centerline{\lsv{10}{A}}
\centerline{\tstrut 24pt\large\lsv{20}{SENTIMENTAL JOURNEY,}}
\centerline{\tstrut 20pt\large \&c.\enspace \&c .}
\vfill
\indent\s{Vol}. II.
\newpage
\null
\newpage
\fontsize{10.5}{15}\selectfont
\parskip 16pt plus 4pt minus 2pt
\pagestyle{folio}
\pagenumbering{arabic}
\thispagestyle{empty}

\hrule
\vskip 4pt
\hrule

\centerline{\tstrut 32pt\lsv{84}{\small THE}}
\stick{\tstrut 32pt \lsv{42}{FILLE}\hss \ls{\scshape De}\hss \lsv{42}{CHAMBRE}}
\head{24pt}{48}{PARIS}

\vskip 6pt

\dropcap{W}{hat} the old French officer\break
had deliver’d upon travel-\break
\stick{ling, bringing Polonius’s advice to}
\stck{his son upon the same subject into my}
\stick{head\tskk and that bringing in Hamlet;}
\stick{and Hamlet, the rest of Shakespear’s}
\stick{works, I stopp’d at the Quai de}
\stick{Conti in my return home, to pur-}
chase the whole set.


\patchv[\baselineskip]{B}{The}

The bookseller said he had not a\break
\stick{set in the world\tskk\i{Comment!} said I;}
\stick{taking one up out of a set which lay}
\stick{upon the counter betwixt us.\tsh He}
\stick{said, they were sent him only to be}
\stick{got bound, and were to be sent back}
\stick{to Versailles in the morning to the}
Count de B\anon.

\tskk And does the Count de B\anon\break
\stick{said I, read Shakespear?  \i{C’est un}}
\stick{\i{Esprit fort}; replied the bookseller.\tskk}
\stick{He loves English books; and what}
\stick{is more to his honour, Monsieur,}
\stick{he love the English too.  You speak}
\stick{this so civilly, said I, that ’tis enough}
\stick{to oblige an Englishman to lay out a}
\stick{Louis d’or or two at your shop\tskk the}
\stick{bookseller made a bow, and was}
\catch{going}
\stck{going to say something, when a young}
\stick{decent girl of about twenty, who by}
\stick{her air and dress, seemed to be \i{fille}}
\stick{\i{de chambre} to some devout woman of}
\stck{fashion, came into the shop and asked}
\stick{for \i{Les Egarments du Cœur \et\ de l’Es}-}
\stick{\i{prit:} the bookseller gave her the book}
\stick{directly; she pulled out a little green}
\stick{sattin purse run round with a ribband}
\stick{of the same colour, and putting her}
\stick{finger and thumb into it, she took out}
\stck{the money, and paid for it. As I had no-}
\stick{thing more to stay me in the shop, we}
both walked out at the door together.

\tsk And what have you to do,\break
\stick{my dear, said I, with \i{The Wanderings}}
\stick{\i{of the Heart}, who scarce know yet}
\stick{you have one? nor till love has}
\catchs{B 2}{first}
\stick{first told you it, or some faithless}
\stck{shepherd has made it ache, can’st thou}
\stick{ever be sure it is so.\tskk \i{Le Dieu m’en}}
\stick{\i{guard!} said the girl.\tskk With reason,}
\stick{said I\tsk for if\sic is a good one, ’tis pity}
\stick{it should be stolen: ’tis a little trea-}
\stick{sure to thee, and gives a better air}
\stick{to your face, than if it was dress’d}
out with pearls.

The young girl listened with a\break
\stick{submissive attention, holding her sat-}
\stick{tin purse by its ribband in her hand}
\stick{all the time\tskk ’Tis a very small one,}
\stick{said I, taking hold of the bottom of}
\stick{it\tskk she held it towards me\tskk and}
\stick{there is very little in it, my dear, said}
\stick{I; but be but as good as thou art hand-}
\stick{some, and heaven will fill it:  I had a}
\catch{parcel}
\stick{parcel of crowns in my hand to pay}
\stick{for Shakespear; and, as she had let go}
\stick{the purse intirely, I put a single one}
\stick{in; and tying up the ribband in a}
bow-knot, returned it to her.

The young girl made me more a\break
\stick{humble courtesy than a low one\tskk}
\stick{’twas one of those quiet, thankful}
\stick{sinkings where the spirit bows itself}
\stick{down\tskk the body does no more than}
\stick{tell it.  I never gave a girl a crown}
\stick{in my life which gave me half the}
pleasure.

My advice, my dear, would not\break
\stick{have been worth a pin to you, said}
\stick{I, if I had not given this along with}
\stick{it: but now, when you see the crown,}
\catchs{B 3}{you’ll}
\stick{you’ll remember it\tskk so don’t, my}
dear, lay it out in ribbands.

Upon my word, Sir, said the\break
\stick{girl, earnestly, I am incapable\tskk in}
\stick{saying which, as is usual in little bar-}
\stick{gains of honour, she gave me her}
\stick{hand\tskk \i{En verite}, \i{Monsieur}, \i{je met}-}
\i{trai cet argent apart}, said she.

When a virtuous convention is\break
\stick{made betwixt man and woman, it}
\stick{sanctifies their most private walks:}
\stick{so notwithstanding it was dusky, yet}
\stick{as both our roads lay the same way,}
\stick{we made no scruple of walking along}
the Quai de Conti together.

\patch[\baselineskip]{She}

She made me a second courtesy in\break
\stick{setting off, and before we got twenty}
\stick{yards from the door, as if she had}
\stick{not done enough before, she made a}
\stick{sort of a little stop to tell me again,}
\tskk she thank’d me.

It was a small tribute, I told her,\break
\stick{which I could not avoid paying to}
\stick{virtue, and would not be mistaken in}
\stick{the person I had been rendering it to}
\stick{for the world\tskk but I see innocence,}
\stick{my dear, in your face\tskk and foul be-}
\stick{fall the man who ever lays a snare in}
its way!

The girl seem’d affected some way\break
\stick{or other with what I said\tskk she gave}
\stick{a low sigh\tskk I found I was not im-}
\catchs{B 4}{powered}
\stick{powered to enquire at all after it\tskk so}
\stick{said nothing more till I got to the}
\stck{corner of the Rue de Nevers, where we}
were to part.

\tskk\tighter{But is this the way, my dear, said I,}\break
\stck{to the hotel de Modene?  she told me it}
\stick{was\tskk or, that I might go by the Rue}
\stick{de Guineygaude, which was the next}
\stick{turn.\tskk Then I’ll go, my dear, by the}
\stick{Rue de Guineygaude, said I, for two}
\stck{reasons; first I shall please myself, and}
\stick{next I shall give you the protection of}
\stick{my company as far on your way as I}
\stick{can.  The girl was sensible I was civil}
\stick{\tskk and said, she wish’d the hotel de}
\stick{Modene was in the Rue de St.\@ Pierre\tsk}
\stick{\tskk You live there? said I.\tskk She told}
\stck{me she was \i{fille de chambre} to Madame}
\catch{R\anon}
\stick{\i{R}\anon\sic\tskk Good God!\@ said I, ’tis the}
\stick{very lady for whom I have brought a}
\stck{letter from Amiens\tskk The girl told me}
\stick{that Madame \i{R}\anon\sic, she believed}
\stick{expected a stranger with a letter, and}
\stk{was impatient to see him\tskk so I desired}
\stick{the girl to present my compliments to}
\stck{Madame R\anon, and say I would cer-}
tainly wait upon her in the morning.

\tighter{We stood still at the corner of the Rue}\break
\stk{de Nevers whilst this pass’d\tskk We then}
\stick{stopp’d a moment whilst she disposed}
\stick{of her \i{Egarments de\sic Cœur}, \&c. more}
\stick{commodiously than carrying them in}
\stick{her hand\tskk they were two volumes;}
\stick{so I held the second for her whilst she}
\stick{put the first into her pocket; and then}
\catch{she}
\stick{she held her pocket, and I put in the}
other after it.

\tight{’Tis sweet to feel by what fine-spun}\break
\stick{threads our affections are drawn to-}
gether.

We set off a-fresh, and as she took\break
\stick{her third step, the girl put her hand}
\stick{within my arm\tskk I was just bidding}
\stick{her\tskk but she did it of herself with}
\stick{that undeliberating simplicity, which}
\stick{shew’d it was out of her head that}
\stick{she had never seen me before.  For}
\stick{my own part, I felt the conviction of}
\stck{consanguinity so strongly, that I could}
\stick{not help turning half round to look in}
\stick{her face, and see if I could trace out}
\catch{any}
\stick{any  thing in it of a family likeness\tsk}
Tut! said I, are we not all relations?

When we arrived at the turning up\break
\stick{of the Rue de Guineygaude, I stopp’d}
\stick{to bid her adieu for good an all: the}
\stick{girl would thank me again for my}
\stick{company and kindness\tskk She bid me}
\stick{adieu twice\tskk I repeated it as often;}
\stck{and so cordial was the parting between}
\stick{us, that had it happen’d any where}
\stick{else, I’m not sure but I should have}
\stick{signed it with a kiss of charity, as}
warm and holy as an apostle.

\tight{But in Paris, as none kiss each other}\break
\stick{but the men\tskk I did, what amounted}
to the same thing\tsh

\tsh I bid God bless her.

\newpage


\head{24pt}{60}{THE PASSPORT}
\head{18pt}{60}{PARIS}

\vskip 2pt


\dropcap{W}{hen} I got home to my\break
hotel, La Fleur told me I\break
\stick{had been enquired after by the Lieu-}
\stick{tenant de Police\tskk The duce take it!}
\stick{said I\tskk\!\tskk I know the reason.  It is}
\stick{time the reader should know it, for}
\stick{in the order of things in which it}
\stick{happened, it was omitted; not that}
\stick{it was out of my head; but that had}
\stick{I told it then, it might have been}
\stick{forgot now\tskk and now is the time I}
want it.

I had left London with so much\break
\stick{precipitation, that it never enter’d my}
\catch{mind}
\stck{mind that we were at war with France;}
\stick{and had reach’d Dover, and look’d}
\stick{through my glass at the hills beyond}
\stick{Boulogne, before the idea presented}
\stick{itself; and with this in its train, that}
\stick{there was no getting there without a}
\stick{passport.  Go but to the end of a}
\stick{street, I have a mortal aversion for}
\stick{returning back no wiser than I sat\sic}
\stick{out; and as this was one of the great-}
\stick{est efforts I had ever made for know-}
\stick{ledge, I could less bear the thoughts}
\stick{of it: so hearing the Count de\anon}
\stick{had hired the packet, I begg’d he}
\stick{would take me in his \i{suite}.  The}
\stick{Count had some little knowledge of}
\stick{me, so made little or no difficulty\tskk}
\stick{only said, his inclination to serve me}
\stick{could reach no farther than Calais;}
\catch{as}
\stck{as he was to return by way of Brussels}
\stick{to Paris: however, when I had once}
\stick{pass’d there, I might get to Paris}
\stick{without interruption; but that in}
\stick{Paris I must make friends and shift}
\stick{for myself.\tskk Let me get to Paris,}
\stick{Monsieur le Count, said I\tskk and I}
\stick{shall do very well.  So I embark’d,}
\stick{and never thought more of the}
matter.

When La Fleur told me the Lieu-\break
\stick{tenant de Police had been enquiring}
\stck{after me\tskk the thing instantly recurred}
\stick{\tskk and by the time La Fleur had well}
\stick{told me, the master of the hotel}
\stick{came into my room to tell me the}
\stick{same thing, with this addition to it,}
\stick{that my passport had been particu-}
\catch{larly}
\stick{larly ask’d after: the master of the}
\stick{hotel concluded with saying, He}
\stick{hoped I had one.\tskk Not I, faith!}
said I.

The master of the hotel retired\break
\stick{three steps from me, as from an in-}
\stick{fected person, as I declared this\tsk and}
\stick{poor La Fleur advanced three steps}
\stick{towards me, and with that sort of}
\stick{movement which a good soul makes}
\stck{to succour a distress’d one\tsk the fellow}
\stick{won my heart by it; and from that}
\stick{single \i{trait}, I knew his character as}
\stick{perfectly, and could rely upon it as}
\stick{firmly, as if he had served me with}
fidelity for seven years.

\patch[\baselineskip]{Mon}

Mon seignior! cried the master of\break
\stick{the hotel\tsk but recollecting himself}
\stick{as he made the exclamation, he in-}
\stick{stantly changed the tone of it\tskk If}
\stick{Monsieur, said he, has not a passport}
\stick{(\i{apparament}) in all likelihood he has}
\stick{friends in Paris who can procure him}
\stick{one.\tskk Not that I know of, quoth I,}
\stick{with an air of indifference.\tskk Then}
\stick{\i{certes}, replied he, you’ll be sent to the}
\stck{Bastile or the Chatelet, \i{au moins}. Poo!}
\stick{said I, the King of France is a good na-}
\stck{tured soul\tskk he’ll hurt no body.\tskk \i{Cela}}
\stick{\i{n’empeche pas}, said he\tskk you will cer-}
\stck{tainly be sent to the Bastile to-morrow}
\stick{morning.\tskk But I’ve taken your lodg-}
\stick{ings for a month, answer’d I, and I’ll}
\stick{not quit them a day before the time}
\stick{for all the kings of France in the}
\catch{world.}
\stick{world.  La Fleur whisper’d in my}
\stick{ear, That no body could oppose the}
king of France.

\i{Pardi!} said my host, \i{ces Messieurs}\break
\stick{\i{Anglois sont des gens tres extraordi}-}
\stick{\i{naires}\tskk and having both said and}
sworn it\tskk he went out.


\patchv[10\baselineskip]{C}{\lss{THE}}


\head{32pt}{64}{THE PASSPORT}
\head{19pt}{0}{The Hotel at Paris}

\vskip 4pt


\initial[nindent=1pt]{I}{COULD} not find in my heart\break
to torture La Fleur’s with a se-\break
\stick{rious look upon the subject of my}
\stick{embarrassment, which was the rea-}
\stick{son I had treated it so cavalierly: and}
\stick{to shew him how light it lay upon}
\stick{my mind, I dropt the subject en-}
\stick{tirely; and whilst he waited upon me}
\stick{at supper, talk’d to him with more}
\stick{than usual gaiety about Paris, and of}
\stick{the opera comique.\tskk La Fleur had}
\stick{been there himself, and had followed}
\stick{me through the streets as far as the}
\stck{bookseller’s shop; but seeing me come}
\catch{out}
\stick{out with the young \i{fille de chambre},}
\stick{and that we walk’d down the Quai de}
\stick{Conti together, La Fleur deem’d it}
\stick{unnecessary to follow me a step fur-}
\stick{ther\tskk so making his own reflections}
\stick{upon it, he took a shorter cut\tsh}
\stick{and got to the hotel in time to be in-}
\stick{form’d of the affair of the Police}
against my arrival.

As soon as the honest creature had\break
\stick{taken away, and gone down to sup}
\stick{himself, I then began to think a}
little seriously about my situation.\tskk

\tskk And here, I know, Eugenius,\break
\stick{thou wilt smile at the remembrance}
\stick{of a short dialogue which pass’d be-}
\catchs{C 2}{twixt}
\stick{twixt us the moment I was going to}
set out\tsh I must tell it here.

Eugenius, knowing that I was as\break
\stick{little subject to be overburthen’d with}
\stick{money as thought, had drawn me}
\stick{aside to interrogate me how much I}
\stick{had taken care for; upon telling him}
\stick{the exact sum, Eugenius shook his}
\stick{head, and said it would not do; so}
\stick{pull’d out his purse in order to empty}
\stick{it into mine.\tskk I’ve enough in con-}
\stick{science, Eugenius, said I.\tsh Indeed,}
\stick{Yorick, you have not, replied Eu-}
\stick{genius\tskk I know France and Italy}
\stick{better than you.\tsh But you don’t}
\stick{consider, Eugenius, said I, refusing}
\stick{his offer, that before I have been}
\stick{three days in Paris, I shall take care}
\catch{to}
\stick{to say or do something or other for}
\stick{which I shall get clapp’d up into the}
\stick{Bastile, and that I shall live there a}
\stick{couple of months entirely at the king}
\stick{of France’s expense.\tskk I beg pardon,}
\stick{said Eugenius, drily: really, I had}
forgot that resource.

Now the event I treated gaily\break
came seriously to my door.

Is it folly, or nonchalance, or phi-\break
\stick{losophy, or pertinacity\tskk or what is}
\stick{it in me, that, after all, when La}
\stick{Fleur had gone down stairs, and I}
\stick{was quite alone, that I could not}
\stick{bring down my mind to think of it}
\stick{otherwise than I had then spoken of}
it to Eugenius?

\patchs{C 3}{\tsk And}


\tskk And as for the Bastile! the ter-\break
\stick{ror is in the word\tskk Make the most}
\stick{of it you can, said I to myself, the}
\stick{Bastile is but another word for a}
\stick{tower\tskk and a tower is but another}
\stick{word for a house you can’t get out of}
\stick{\tskk Mercy on the gouty!  for they are}
\stick{in it twice a year\tskk But with nine}
\stick{livres a day, and pen and ink and}
\stick{paper and patience, albeit a man can’t}
\stick{get out, he may do very well within}
\stck{\tskk at least for a month or six weeks; at}
\stick{the end of which, if he is a harmless}
\stick{fellow his innocence appears, and}
\stick{he comes out a better and wiser man}
than he went in.

I had some occasion (I forget what)\break
\stick{to step into the court-yard, as I settled}
\catch{this}
\stick{this account; and remember I walk’d}
\stick{down stairs in no small triumph with}
\stick{the conceit of my reasoning.\tskk Be-}
\stck{shrew the \i{sombre} pencil! said I vaunt-}
\stick{ingly\tskk for I envy not its powers,}
\stick{which paints the evils of life with so}
\stick{hard and deadly a colouring.  The}
\stick{mind sits terrified at the objects she}
\stck{has magnified herself, and blackened:}
\stick{reduce them to their proper size and}
\stick{hue she overlooks them\tskk ’Tis true,}
\stick{said I, correcting the proposition\tskk}
\stick{the Bastile is not an evil to be de-}
\stck{spised\tskk but strip it of its towers\tskk fill}
\stk{up the fossè\tskk unbarricade the doors\tskk}
\stick{call it simply a confinement, and sup-}
\stck{pose ’tis some tyrant of a distemper\tskk}
\stick{and not of a man which holds you in}
\catchs{C 4}{it\tskk}
\stick{it\tskk the evil vanishes, and you bear}
the other half without complaint.

I was interrupted in the hey-day\break
\stick{of this soliloquy, with a voice which}
\stick{I took to be of a child, which com-}
\stick{plained \lqq it could not get out.\rqq \tskk I}
\stick{look’d up and down the passage, and}
\stick{seeing neither man, woman, or child,}
I went out without farther attention.

In my return back through the\break
\stick{passage, I heard the same words re-}
\stick{peated twice over; and looking up,}
\stick{I saw it was a starling hung in a little}
\stick{cage.\tskk \lqq I can’t get out\tskk I can’t get}
\lqq out,\rqq\ said the starling.

\patch[\baselineskip]{I stood}

I stood looking at the bird: and\break
\stck{to every person who came through the}
\stick{passage it ran fluttering to the side to-}
\stick{wards which they approach’d it, with}
\stk{the same lamentation of its captivity\tskk}
\stick{\lqq I can’t get out\rqq, said the starling\tskk}
\stick{God help thee! said I, but I’ll let}
\stick{thee out, cost what it will; so I turn’d}
\stick{about the cage to get to the door; it}
\stck{was twisted and double twisted so fast}
\stick{with wire, there was no getting it}
\stick{open without pulling the cage to}
pieces\tskk I took both hands to it.

The bird flew to the place where\break
\stick{I was attempting his deliverance, and}
\stick{thrusting his head through the trellis,}
\stick{press’d his breast against it, as if im-}
\stick{patient\tskk I fear, poor creature!\@ said I,}
\catch{I cannot}
\stick{I cannot set thee at liberty\tskk \lqq No,\rqq}
\stick{said the starling\tskk \lqq I can’t get out}
\stick{\lqq\tskk I can’t get out,\rqq\ said the star-}
ling.

I vow, I never had my affections\break
\stick{more tenderly awakened; or do I re-}
\stick{member an incident in my life, where}
\stick{the dissipated spirits, to which my}
\stick{reason had been a bubble, were so}
\stick{suddenly call’d home. Mechanical}
\stck{as the notes were, yet so true in tune to}
\stick{nature were they chanted, that in one}
\stick{moment they overthrew all my syste-}
\stick{matic reasonings upon the Bastile;}
\stick{and I heavily walk’d up stairs, un-}
\stick{saying every word I had said in go-}
ing down them.

\patch[.66\baselineskip]{Disguise}

Disguise thyself as thou wilt, still\break
\stick{slavery! said I\tskk still thou art a bit}
\stick{ter draught; and though thousands}
\stick{in all ages have been made to drink}
\stick{of thee, thou art no less bitter on}
\stck{that account.\tskk ’tis thou, thrice sweet}
\stck{and gracious goddess, addressing my-}
\stick{self to \s{Liberty}, whom all in pub-}
\stick{lic or in private worship, whose taste}
\stick{is grateful, and ever wilt be so, till}
\stk{\s{Nature} herself shall change\tskk No \i{tint}}
\stick{of words can spot thy snowy mantle,}
\stick{or chymic power turn thy sceptre}
\stick{into iron\tskk with thee to smile upon}
\stick{him as he eats his crust, the swain}
\stick{is happier than his monarch, from}
\stk{whose court thou art exiled\tskk Gracious}
\stick{heaven!\@ cried I, kneeling down upon}
\stick{the last step but one in my ascent\tsk}
\catch{grant}
\stick{grant me but health, thou great Be-}
\stick{stower of it, and give me but this}
\stick{fair goddess as my companion\tskk}
\stick{and shower down thy mitres, if it}
\stick{seems good unto thy divine provi-}
\stick{dence, upon those heads which are}
aching for them.

\patch[10\baselineskip]{\lsss{THE}}


\head{28pt}{48}{THE CAPTIVE}
\head{18pt}{64}{PARIS}

\vskip 2pt

\dropcap{T}{he} bird in his cage pursued\break
me into my room; I sat down\break
\stick{close to my table, and leaning my}
\stick{head upon my hand, I began to}
\stick{figure to myself the miseries of con-}
\stick{finement.  I was in a right frame}
\stick{for it, and so I gave full scope to my}
imagination.

I was going to begin with the\break
\stick{millions of my fellow creatures born}
\stick{to no inheritance but slavery; but}
\stick{finding, however affecting the pic-}
\stick{ture was, that I could not bring it}
\catch{near}
\stick{near me, and that the multitude of}
\stick{of\sic sad groups in it did but distract}
me.\tskk

\tskk I took a single captive, and hav-\break
\stick{ing first shut him up in his dungeon,}
\stick{I then look’d through the twilight of}
his grated door to take his picture.

I beheld his body half wasted away\break
\stick{with long expectation and confine-}
\stck{ment, and felt what kind of sickness of}
\stck{the heart it was which arises from hope}
\stick{deferr’d.  Upon looking nearer I saw}
\stick{him pale and feverish: in thirty years}
\stck{the western breeze had not once fann’d}
\stick{his blood\tskk he had seen no sun, no}
\stick{moon in all that time\tskk nor had the}
\catch{voice}
\stick{voice of friend or kinsman breathed}
through his lattice\tskk his children\tskk

\tsk But here my heart began to bleed\break
\stick{\tskk and I was forced to go on with}
another part of the portrait.

He was sitting upon the ground\break
\stick{upon a little straw, in the furthest}
\stick{corner of his dungeon, which was}
\stick{alternately his chair and bed: a little}
\stck{calendar of small sticks were laid at the}
\stick{head notch’d all over with the dismal}
\stck{days and nights he had pass’d there\tskk}
\stick{he had one of these little sticks in his}
\stick{hand, and with a rusty nail he was}
\stick{etching another day of misery to add}
\stick{to the heap.  As I darkened the little}
\stick{light he had, he lifted up a hopeless}
\catch{eye}
\stick{eye towards the door, then cast it}
\stick{down\tskk shook his head, and went on}
\stick{with his work of affliction.  I heard}
\stick{his chains upon his legs, as he turn’d}
\stick{his body to lay his little stick upon the}
\stck{bundle\tskk He gave a deep sigh\tskk I saw}
\stick{the iron enter into his soul\tskk I burst}
\stick{into tears\tskk I could not sustain the}
\stick{picture of confinement which my}
\stick{fancy had drawn\tskk I started up from}
\stick{my chair, and calling La Fleur, I bid}
\stick{him bespeak me a \i{remise}, and have}
\stick{it ready at the door of the hotel by}
nine in the morning.

\tsk I’ll go directly, said I, myself to\break
Monsieur le Duc de Choiseul.

\patch[\baselineskip]{La}

La Fleur would have put me to\break
\stick{bed; but not willing he should see}
\stick{any thing upon my cheek, which}
\stick{would cost the honest fellow a heart}
\stick{ache\tskk I told him I would go to bed}
\stick{by himself\sic\tskk and bid him go do the}
same.


\patchv[12\baselineskip]{D}{\lsss{THE}}
\thispagestyle{minusone}
\head{32pt}{48}{THE STARLING}
\head{24pt}{28}{ROAD {\footnotesize TO} VERSAILLES}

\vskip 4pt


\initial{I}{\ls{GOT}} into my \i{remise} the hour I\break
proposed: La Fleur got up\break
\stick{behind, and I bid the coachman}
\stick{make the best of his way to Ver-}
sailles.

As there was nothing in this road,\break
\stick{or rather nothing which I look for}
\stick{in travelling, I cannot fill up the}
\stick{blank better than with a short history}
\stick{of this self-same bird, which became}
the subject of the last chapter.


\patch[\baselineskip]{Whilst}
\thispagestyle{minusone}

Whilst the Honourable Mr.\@ \anon\break
\stick{was waiting for a wind at Dover}
\stick{it had been caught upon the cliffs be-}
\stick{fore it could well fly, by an English}
\stick{lad who was his groom; who not}
\stick{caring to destroy it, had taken it in}
\stick{his breast into the packet\tskk and by}
\stick{course of feeding it, and taking it}
\stick{once under his protection, in a day}
\stick{or two grew fond of it, and got it}
safe along with him to Paris.

At Paris the lad had laid out a\sic\break
\stick{a livre in a little cage for the starling,}
\stick{and as he had little to do better the}
\stck{five months his master stay’d there, he}
\stick{taught it in his mother’s tongue the}
\stck{four simple words\tskk (and no more)\tskk}
\catchs{D 2}{to}
\stick{to which I own’d myself so much}
it’s\sic debtor.

Upon his master’s going on for\break
\stick{Italy\tsk the lad had given it to the}
\stick{master of the hotel\tsk But his little}
\stick{song for liberty, being in an \i{unknown}}
\stick{language at Paris\tsk the bird had little}
\stick{or no store set by him\tsk so La Fleur}
\stick{bought both him and his cage for}
me for a bottle of Burgundy.

In my return from Italy I brought\break
\stick{him with me to the country in whose}
\stick{language he had learn’d his notes\tsk}
\stick{and telling the story of him to Lord}
\stick{A\tskk Lord A begg’d the bird of me\tskk}
\stick{in a week Lord A gave him to}
\stick{Lord B\tskk Lord B made a present}
\stick{of him to Lord C\tskk and Lord’s\sic C’s}
\catch{gentle-}
\stick{gentleman sold him to Lord D’s for}
\stick{a shilling\tskk Lord D gave him to}
\stick{Lord E\tskk and so on\tskk half round the}
\stick{alphabet\tskk  From that rank he pass’d}
\stick{into the lower house, and pass’d the}
\stick{hands of as many commoners\tsh}
\stick{But as all these wanted to \i{get in}\tsk and}
\stick{my bird wanted to get out\tsk he had}
\stick{almost as little store set by him in}
London as in Paris.

It is impossible but many of my\break
\stick{readers must have heard of him; and}
\stick{if any by mere chance have ever seen}
\stick{him\tskk I beg leave to inform them,}
\stick{that that bird was my bird\tskk or some}
vile copy set up to represent him.

I have nothing further to add upon\break
\stick{him, but that from that time to this,}
\catchs{D 3}{I have}
\stick{I have borne this poor starling as the}
crest to my arms.\tskk Thus:\vrule width 0pt depth 14pt
$$\includegraphics[width=\textwidth]{the-yorick-arms}$$
\indent\tsh And let the heralds officers\break
twist his neck about if they dare.

\patch[.67\baselineskip]{\lss{THE}}


\head{34pt}{60}{THE\enspace ADDRESS}
\head{22pt}{44}{VERSAILLES}

\vskip 6pt


\initial{I}{\ls{SHOULD}} not like to have my\break
enemy take a view of my mind,\break
\stick{when I am going to ask protection}
\stick{of any man: for which reason I ge-}
\stick{nerally endeavour to protect myself;}
\stick{but this going to Monsieur Le Duc}
\stick{de C\thinspace\astv\ was an act of compul-}
\stick{sion\tsk had it been an act of choice,}
\stick{I should have done it, I suppose, like}
other people.

How many mean plans of dirty\break
\stick{address, as I went along, did my}
\catchs{D 4}{servile}
\stick{servile heart form!  I deserved the}
Bastile for every one of them.

Then nothing would serve me,\break
\stick{when I got within sight of Versailles,}
\stck{but putting words and sentences toge-}
\stick{ther, and conceiving attitudes and}
\stick{tones to wreath myself into Monsieur}
\stick{Le Duc de C\thinspace\astv’s good graces\tskk}
\stick{This will do\tsh said I\tsk Just as}
\stick{well, retorted I again, as a coat}
\stick{carried up to him by an adventurous}
\stick{taylor, without taking his measure\tsk}
\stick{Fool!\@ continued I\tskk see Monsieur Le}
\stick{Duc’s face first\tskk observe what cha-}
\stick{racter is written in it; take notice in}
\stick{what posture he stands to hear you}
\stick{\tskk mark the turns and expressions of}
\catch{his}
\stick{his body and limbs\tskk And for the tone}
\stick{\tskk the first sound which comes from}
\stick{his lips will give it you; and from}
\stick{all these together you’ll compound}
\stick{an address at once upon the spot,}
\stick{which cannot disgust the Duke\tskk the}
\stick{ingredients are his own, and most}
likely to go down.

Well!\@ said I, I wish it well over\tskk\break
\stick{Coward again! as if man to man was}
\stick{not equal throughout the whole sur-}
\stick{face of the globe; and if in the field}
\stick{\tskk why not face to face in the cabinet}
\stick{too?  And trust me, Yorick, when-}
\stick{ever it is not so, man is false to him-}
\stick{self; and betrays his own succours}
\stick{ten times, where nature does it once.}
\catch{Go}
\stick{Go to the Duc de C\anon\ with the}
\stick{Bastile in thy looks\tskk My life for it,}
\stick{thou wilt be sent back to Paris in half}
an hour, with an escort.

I believe so, said I\tskk Then I’ll go\break
\stick{to the Duke, by heaven! with all the}
gaity\sic and debonairness in the world.\tskk

\tskk And there you are wrong again,\break
\stick{replied I\tskk A heart at ease, Yorick,}
\stick{flies into no extremes\tskk ’tis ever on}
\stick{its center.\tskk Well! well! cried I, as}
\stick{the coachman turn’d in at the gates}
\stick{\tsk I find I shall do very well: and by}
\stick{the time he had wheel’d round the}
\stick{court, and brought me up to the}
\stick{door, I found myself so much the}
\stick{better for my own lecture, that I}
\catch{neither}
\stick{neither ascended the steps like a}
\stick{victim to justice, who was to part}
\stick{with life upon the topmost,\tskk nor}
\stick{did I mount them with a skip and a}
\stick{couple of strides, as I do when I}
fly up, Eliza!\@ to thee, to meet it.

As I enter’d the door of the saloon,\break
\stick{I was met by a person who possibly}
\stick{might be the maitre d’hotel, but had}
\stick{more the air of one of the under}
\stick{secretaries, who told me the Duc de}
\stick{C\anon\ was busy\tskk I am utterly ig-}
\stick{norant, said I, of the forms of ob-}
\stick{taining an audience, being an absolute}
\stick{stranger, and what is worse in the pre-}
\stick{sent conjuncture of affairs, being an}
\stick{Englishman too.\tsh He replied, that}
\stick{did not increase the difficulty.\tskk I made}
\catch{him}
\stick{him a slight bow, and told him, I}
\stick{had something of importance to say}
\stick{to Monsieur Le Duc.  The secretary}
\stick{look’d towards the stairs, as if he}
\stick{was about to leave me to carry up}
\stick{this account to some one\tskk But I}
\stick{must not mislead you, said I\tskk  for}
\stick{what I have to say is of no manner}
\stick{of importance to Monsieur Le Duc}
\stick{de C\anon\tskk but of great import-}
\stick{ance to myself.\tskk \i{C’est une autre affaire},}
\stick{replied he\tshh Not at all, said I,}
\stick{to a man of gallantry.\tskk But pray,}
\stick{good sir, continued I, when can a}
\stick{stranger hope to have \i{accesse?} In not}
\stick{less than two hours, said he, looking}
\stick{at his watch.  The number of equi-}
\stick{pages in the court-yard seem’d to}
\stick{justify the calculation, that I could}
\catch{have}
\stick{have no nearer a prospect\tskk and as}
\stick{walking backwards and forwards in}
\stick{the saloon, without a soul to commune}
\stick{with, was for the time as bad as}
\stick{being in the Bastile itself, I instantly}
\stick{went back to my \i{remise}, and bid the}
\stick{coachman drive me to the \i{cordon bleu},}
which was the nearest hotel.

I think there is a fatality in it\tskk I\break
seldom go to the place I set out for.

\patch[6\baselineskip]{\lss{LE}}



\head{32pt}{64}{LE\enspace PATISSIER}
\head{18pt}{44}{VERSAILLES}

\vskip 6pt

\dropcap{B}{efore} I had got half-way\break
down the street I changed my\break
\stick{mind: as I am at Versailles, thought}
\stick{I, I might as well take a view of the}
\stick{town; so I pull’d the cord, and or-}
\stick{dered the coachman to drive round}
\stick{some of the principal streets\tskk I sup-}
\stick{pose the town is not very large, said}
\stick{I.\tskk The coachman begg’d pardon for}
\stick{setting me right, and told me it was}
\stick{very superb, and that numbers of the}
\stick{first dukes and marquises and counts}
\stick{had hotels\tskk The Count de B\anon, of}
\catch{whom}
\stick{whom the bookseller at the Quai de}
\stick{Conti had spoke so handsomely the}
\stick{night before, came instantly into my}
\stick{mind.\tskk And why should I not go,}
\stick{thought I, to the Count de B\anon,}
\stick{who has so high an idea of English}
\stick{books and Englishmen\tskk and tell}
\stick{him my story?\@ so I changed my}
\stick{mind a second time\tskk In truth it was}
\stick{the third; for I had intended that day}
\stick{for Madame de R\anon\ in the Rue}
\stick{St.\@ Pierre, and had devoutly sent her}
\stick{word by her \i{fille de chambre} that I}
\stick{would assuredly wait upon her\tskk but}
\stick{I am govern’d by circumstances\tskk I}
\stick{cannot govern them: so seeing a man}
\stick{standing with a basket on the other}
\stick{side of the street, as if he had some-}
\stick{thing to sell, I bid La Fleur go up}
\catch{to}
\stick{to him, and enquire for the Count’s}
hotel.

La Fleur return’d a little pale; and\break
\stick{told me it was a Chevalier de St.}
\stick{Louis selling \i{patès}\tskk It is impossible,}
\stick{La Fleur!\@ said I.\tskk La Fleur could no}
\stick{more account for the phenomenon}
\stick{than myself; but persisted in his}
\stick{story: he had seen the croix set in}
\stick{gold, with its red ribband, he said, tied}
\stick{to his button-hole\tskk and had look’d}
\stick{into the basket and seen the \i{patès}}
\stick{which the Chevalier was selling; so}
could not be mistaken in that.

Such a reverse in man’s life awakens\break
\stick{a better principle than curiosity: I}
\stick{could not help looking for some time}
\catch{at}
\stick{at him as I sat in the \i{remise}\tskk the}
\stick{more I look’d at him\tsk his croix and}
\stick{his basket, the stronger they wove}
\stick{themselves into my brain\tskk I got out}
of the \i{remise} and went towards him.

He was begirt with a clean linen\break
\stick{apron which fell below his knees, and}
\stick{with a sort of a bib that went half way}
\stick{up his breast; upon the top of this,}
\stick{but a little below the hem, hung his}
\stick{croix.  His basket of little \i{patès} was}
\stick{cover’d over with a white damask}
\stick{napkin; another of the same kind}
\stick{was spread at the bottom; and there}
\stick{was a look of \i{propreté} and neatness}
\stick{throughout; that one might have}
\stick{bought his \i{patès} of him, as much}
from appetite as sentiment.

\patchv{E}{He}

He made an offer of them to nei-\break
\stick{ther; but stood still with them at the}
\stick{corner of a hotel, for those to buy}
who chose it, without solicitation.

He was about forty-eight\tskk of a\break
\stick{sedate look, something approaching}
\stick{to gravity.  I did not wonder.\tskk I}
\stick{went up rather to the basket than}
\stick{him, and having lifted up the nap-}
\stick{kin, and taking one of his \i{patès} into}
\stick{my hand\tskk I begg’d he would explain}
the appearance which affected me.

He told me in a few words, that\break
\stick{the best part of his life had pass’d in}
\stick{the service, in which, after spending}
\stick{a small patrimony, he had obtain’d}
\stick{a company and the croix with it;}
\catch{but}
\stick{but that, at the conclusion of the last}
\stick{peace, his regiment being reformed,}
\stick{and the whole corps, with those of}
\stick{some other regiments, left without any}
\stick{provision\tskk he found himself in a}
\stick{wide world without friends, without}
\stick{a livre\tskk and indeed, said he, with-}
\stick{out any thing but this\tskk (pointing, as}
\stick{he said it, to his croix)\tskk The poor}
\stick{chevalier won my pity, and he finish’d}
\stick{the scene with winning my esteem}
too.

The king, he said, was the most\break
\stick{generous of princes, but his gene-}
\stick{rosity could neither relieve nor reward}
\stick{every one, and it was only his mis-}
\stick{fortune to be amongst the number.}
\stick{He had a little wife, he said, whom}
\catchs{E 2}{he}
\stick{he loved, who did the \i{patisserie}; and}
\stick{added, he felt no dishonour in defend-}
\stick{ing her and himself from want in this}
\stick{way\tskk unless Providence had offer’d}
him a better.

It would be wicked to with-hold a\break
\stick{pleasure from the good, in passing}
\stick{over what happen’d to this poor}
\stick{Chevalier of St.\@ Louis about nine}
months after.

It seems he usually took his stand\break
\stick{near the iron gates which lead up to}
\stick{the palace, and as his croix had}
\stick{caught the eyes of numbers, num-}
\stick{bers had made the same enquiry}
\stick{which I had done\tskk He had told}
\stick{them the same story, and always with}
\catch{so}
\stick{so much modesty and good sense,}
\stick{that it had reach’d at last the king’s}
\stick{ears\tskk who, hearing the Chevalier}
\stick{had been a gallant officer, and re-}
\stick{spected by the whole regiment as a}
\stick{man of honour and integrity\tskk he}
\stick{broke up his little trade by a pension}
of fifteen hundred livres a year.

As I have told this to please the\break
\stick{reader, I beg he will allow me to}
\stick{relate another out of its order, to}
\stick{please myself\tskk the two stories reflect}
\stick{light upon each other,\tskk and ’tis a pity}
they should be parted.

\patchs[3\baselineskip]{E 3}{\lss{THE}}

\head{32pt}{64}{THE\enspace SWORD}
\head{24pt}{48}{RENNES}

\vskip 6pt

\dropcap{W}{hen} states and empires have\break
their periods of declension,\break
\stick{and feel in their turns what distress}
\stick{and poverty is\tskk I stop not to tell}
\stick{the causes which gradually brought}
\stick{the house d’E\anon\ in Britany}
\stick{into decay.  The Marquis d’E\anon}
\stick{had fought up against his condition}
\stick{with great firmness; wishing to pre-}
\stick{serve, and still shew to the world,}
\stick{some little fragments of what his an-}
\stick{cestors had been\tskk their indiscretions}
\stick{had put it out of his power. There}
\catch{was}
\stick{was enough left for the little exigen-}
\stick{cies of \i{obscurity}\tskk But he had two boys}
\stick{who looked up to him for \i{light}\tskk he}
\stick{thought they deserved it.  He had}
\stick{tried his sword\tskk it could not open}
\stick{the way\tskk the \i{mounting} was too ex-}
\stick{pensive\tskk and simple œconomy was}
\stick{not a match for it\tskk there was no re-}
source but commerce.

In any other province in France,\break
\stick{save Britany, this was smiting the}
\stick{root for ever of the little tree his}
\stick{pride and affection wish’d to see re-}
\stick{blossom\tskk But in Britany, there being}
\stick{a provision for this, he avail’d himself}
\stick{of it; and taking an occasion when}
\stick{the states were assembled at Rennes,}
\catchs{E 4}{the}
\stick{the Marquis, attended with his two}
\stick{boys, enter’d the court; and having}
\stick{pleaded the right of an ancient law}
\stick{of the duchy, which, though seldom}
\stick{claim’d, he said, was no less in force;}
\stick{he took his sword from his side\tskk Here}
\stick{\tskk said he\tskk take it; and be trusty}
\stick{guardians of it, till better times put}
me in condition to reclaim it.

The president accepted the Mar-\break
\stick{quis’s sword\tskk he stay’d a few minutes}
\stick{to see it deposited in the archives of}
his house\tskk and departed.

The Marquis and his whole family\break
\stick{embarked the next day for Marti-}
\stick{nico, and in about nineteen or twenty}
\stick{years of successful application to bu-}
\catch{siness,}
\stick{siness, with some unlook’d for be-}
\stick{quests from distant branches of his}
\stick{house\tskk return’d home to reclaim his}
nobility and to support it.

It was an incident of good for-\break
\stick{tune which will never happen to any}
\stick{traveller, but a sentimental one, that}
\stick{I should be at Rennes at the very}
\stick{time of this solemn requisition: I}
call it solemn\tskk it was so to me.

The Marquis enter’d the court with\break
\stick{his whole family: he supported his}
\stick{lady\tskk his eldest son supported his}
\stck{sister, and his youngest was at the other}
\stick{extreme of the line next his mother.}
\stick{\tskk he put his handkerchief to his}
face twice\tskk

\patch{\tskk There}

\tskk There was a dead silence.  When\break
\stick{the Marquis had approach’d within}
\stick{six paces of the tribunal, he gave the}
\stick{Marchioness to his youngest son, and}
\stick{advancing three steps before his fa-}
\stick{mily\tskk he reclaim’d his sword.  His}
\stick{sword was given him, and the moment}
\stick{he got it into his hand he drew it al-}
\stick{most out of the scabbard\tskk ’twas the}
\stick{shining face of a friend he had once}
\stick{given up\tskk he look’d attentively along}
\stick{it, beginning at the hilt, as if to see}
\stick{whether it was the same\tskk when, ob-}
\stick{serving a little rust which it had con-}
\stick{tracted near the point, he brought it}
\stick{near his eye, and bending his head}
\stick{down over it\tskk I think I saw a tear}
\stick{fall upon the place:  I could not be}
deceived by what followed.

\patch{\lqq I shall}

\lqq I shall find, said he, some \i{other}\break
\lqq \i{way}, to get it off.\rqq

When the Marquis had said this,\break
\stick{he return’d his sword into its scabbard,}
\stick{made a bow to the guardians of it\tskk}
\stick{and, with his wife and daughter and}
\stick{his two sons following him, walk’d}
out.

O how I envied him his feelings!


\patch[6\baselineskip]{\lsss{THE}}

\head{34pt}{68}{THE\enspace PASSPORT}
\head{22pt}{44}{VERSAILLES}
\vskip 4pt

\initial{I}{FOUND} no difficulty in get-\break
ting admittance to Monsieur le\break
\stick{Count de B\anon.  The set of Shake-}
\stick{spears was laid upon the table, and}
\stick{he was tumbling them over.  I}
\stick{walk’d up close to the table, and}
\stick{giving first such a look at the books}
\stick{as to make him conceive I knew}
\stick{what they were\tskk I told him I had}
\stick{come without any one to present me,}
\stick{knowing I should meet with a friend}
\stick{in his apartment, who, I trusted, would}
\stick{do it for me\tskk it is my countryman}
\catch{the}
\stick{the great Shakespear, said I, point-}
\stick{ing to his works\tskk \i{et ayez la bontè},}
\stick{\i{mon cher ami}, apostrophizing his}
\stick{spirit, added I, \i{de me faire cet hon}-}
\i{neur-la}.\tsh

The Count smil'd at the singu-\break
\stick{larity of the introduction; and see-}
\stick{ing I look’d a little pale and sickly,}
\stick{insisted upon my taking an arm-chair:}
\stick{so I sat down; and to save him con-}
\stick{jectures upon a visit so out of all}
\stick{rule, I told him simply of the inci-}
\stick{dent in the bookseller’s shop, and}
\stick{how that had impell’d me rather}
\stick{to go to him with the story of a little}
\stick{embarrassment I was under, than to}
\stick{any other man in France\tskk And what}
\catch{is}
\stick{is your embarrassment? let me hear}
\stick{it, said the Count.  So I told}
\stick{him the story just as I have told it}
the reader\tsk

\tskk And the master of my hotel,\break
\stick{said I, as I concluded it, will needs}
\stick{have it, Monsieur le Count, that I}
\stick{shall be sent to the Bastile\tskk but I}
\stick{have no apprehensions, continued I\tskk}
\stick{for in falling into the hands of the}
\stick{most polish’d people in the world,}
\stick{and being conscious I was a true man,}
\stick{and not come to spy the nakedness of}
\stick{the land, I scarce thought I laid at}
\stick{their mercy.\tskk It does not suit the gal-}
\stick{lantry of the French, Monsieur le}
\stick{Count, said I, to shew it against in-}
valids.

\patch{An}

An animated blush came into the\break
\stick{Count de B\anon’s cheeks, as I}
\stick{spoke this\tskk \i{Ne craignez rien}\tskk Don’t}
\stick{fear, said he\tskk Indeed, I don’t, replied I}
\stick{again\tskk besides, continued I a little}
\stick{sportingly\tskk I have come laughing all}
\stick{the way from London to Paris, and}
\stick{I do not think Monsieur le Duc de}
\stick{Choiseul is such an enemy to mirth,}
\stick{as to send me back crying for my}
pains.

\tshh My application to you, Mon-\break
\stick{sieur le Count de B\anon\ (mak-}
\stick{ing him a low bow) is to desire he}
will not.

The Count heard me with great\break
\stick{good nature, or I had not said half}
\catch{as}
\stick{as much\tskk and once or twice said\tskk}
\stick{\i{C’est bien dit}.  So I rested my cause}
\stick{there\tskk and determined to say no}
more about it.

The Count led the discourse: we\break
\stick{talk’d of indifferent things;\tskk of books}
\stick{and politicks, and men\tskk and then of}
\stick{women\tskk God bless them all! said I,}
\stick{after much discourse about them\tskk}
\stick{there is not a man upon earth who}
\stick{loves them so much as I do: after}
\stick{all the foibles I have seen, and all the}
\stick{satires I have read against them, still}
\stick{I love them; being firmly persuaded}
\stick{that a man who has not a sort of an}
\stick{affection for the whole sex, is incapa-}
\stick{ble of ever loving a single one as he}
ought.

\patch{\i{Heh}}

\i{Hèh bien! Monsieur l’Anglois}, said\break
\stick{the Count, gaily \tsk\ You are not}
\stick{come to spy the nakedness of the}
\stick{land\tskk I believe you\tskk \i{ni encore},}
\stick{I dare say, \i{that} of our women\tskk But}
\stick{permit me to conjecture\tskk if, \i{par}}
\stick{\i{hazard}, they fell in your way\tskk that}
the prospect would not affect you.

I have something within me which\break
\stick{cannot bear the shock of the least in-}
\stick{decent insinuation: in the sportability}
\stick{of chit-chat I have often endeavoured}
\stick{to conquer it, and with infinite pain}
\stick{have hazarded a thousand things to a}
\stick{dozen of the sex together\tskk the least}
\stick{of which I could not venture to a}
single one, to gain heaven.

\patchv[\baselineskip]{F}{Excuse}

Excuse me, Monsieur Le Count,\break
\stick{said I\tskk as for the nakedness of your}
\stick{land, if I saw it, I should cast my}
\stick{eyes over it with tears in them\tskk and}
\stick{for that of your women (blushing at}
\stick{the idea he had excited in me) I am}
\stick{so evangelical in this, and have such}
\stick{a fellow-feeling for what ever is \i{weak}}
\stick{about them, that I would cover it}
\stick{with a garment, if I knew how to}
\stick{throw it on\tskk But I could wish, con-}
\stick{tinued I, to spy the \i{nakedness} of}
\stick{their hearts, and through the different}
\stick{disguises of customs, climates, and}
\stick{religion, find out what is good in}
\stick{them to fashion my own by\tskk and}
therefore am I come.

\patch{It}

It is for this reason, Monsieur le\break
\stick{Compte, continued I, that I have}
\stick{not seen the Palais royal\tskk nor the}
\stick{Luxembourg\tskk nor the Façade of}
\stick{the Louvre\tskk nor have attempted to}
\stick{swell the catalogues we have of pic-}
\stick{tures, statues, and churches\tskk I con-}
\stick{ceive every fair being as a temple,}
\stick{and would rather enter in, and see the}
\stick{original drawings and loose sketches}
\stick{hung up in it, than the transfigura-}
tion of Raphael itself.

The thirst of this, continued I,\break
\stick{as impatient as that which inflames}
\stick{the breast of the connoisseur, has}
\stick{led me from my own home into}
\stick{France\tskk and from France will lead}
\stick{me through Italy\tskk ’tis a quiet journey}
\catchs{F 2}{of}
\stick{of the heart in pursuit of \s{Nature},}
\stick{and those affections which rise out of}
\stick{her, which make us love each other}
\tskk and the world, better than we do.

The Count said a great many civil\break
\stick{things to me upon the occasion; and}
\stick{added very politely how much he}
\stick{stood obliged to Shakespear for}
\stick{making me known to him\tskk but,}
\stick{\i{a-propos}, said he\tskk Shakespear is full}
\stick{of great things\tskk He forgot a small}
\stick{punctilio of announcing your name}
\stick{\tskk it puts you under a necessity of}
doing it yourself.

\patch[3\baselineskip]{\lss{THE}}

\head{34pt}{68}{THE\enspace PASSPORT}
\head{22pt}{44}{VERSAILLES}
% \vskip 4pt

\dropcap{T}{here} is not a more per-\break
plexing affair in life to me,\break
\stick{than to set about telling any one who}
\stick{I am\tskk for there is scarce any body I}
\stick{cannot give a better account of than}
\stick{of myself; and I have often wish’d I}
\stick{could do it in a single word\tskk and}
\stick{have an end of it.  It was the only}
\stick{time and occasion in my life, I could}
\stick{accomplish this to any purpose\tskk for}
\stick{Shakespear lying upon the table, and}
\stick{recollecting I was in his books, I}
\stick{took up Hamlet, and turning imme-}
\stick{diately to the grave-diggers scene in}
\catch{the\sic}
\stick{the fifth act, I lay’d my finger upon}
\stick{\s{Yorick}, and advancing the book to}
\stick{the Count, with my finger all the}
\stick{way over the name \tskk\ Me, \i{Voici!}}
said I.

Now, whether the idea of poor\break
\stick{Yorick’s skull was put out of the}
\stick{Count’s mind, by the reality of my}
\stick{own, or by what magic he could}
\stick{drop a period of seven or eight hun-}
\stick{dred years, makes nothing in this}
\stick{account\tskk ’tis certain the French con-}
\stick{ceive better than they combine\tskk I}
\stick{wonder at nothing in this world, and}
\stick{the less at this; inasmuch as one of}
\stick{the first of our own church, for whose}
\stick{candour and paternal sentiments I}
\stick{have the highest veneration, fell}
\catch{into}
\stick{into the same mistake in the very}
\stick{same case.\tskk \lqq He could not bear, he}
\stick{\lqq\ said, to look into sermons wrote}
\stick{\lqq\ by the king of Denmark’s jester.\rqq}
\stick{\tskk Good, my lord!\@ said I\tskk but there}
\stick{are two Yorick’s.\sic  The Yorick your}
\stick{lordship thinks of, has been dead}
\stick{and buried eight hundred years ago;}
\stick{he flourish’d in Horwendillus’s court\tsk}
\stick{the other Yorick is myself, who have}
\stick{flourish’d my lord in no court\tskk He}
\stick{shook his head\tskk Good God! said I,}
\stick{you might as well confound Alex-}
\stick{ander the Great, with Alexander the}
\stick{Copper-smith, my lord\tsh ’Twas all}
one, he replied\tskk

\tskk If Alexander king of Macedon\break
\stick{could have translated your lordship,}
\catchs{F 4}{said}
\stick{said I \tskk\ I’m sure your Lordship}
would not have said so.

The poor Count de B\anon\ fell\break
but into the same \i{error}\tsk

\tshh \i{Et, Monsieur, est il Yorick?}\break
\stick{cried the Count.\tskk \i{Je le suis}, said I.}
\stick{\tskk \i{Vous?}\tskk \i{Moi}\tskk \i{moi qui ai l’honneur}}
\stick{\i{de vous parler, Monsieur le Compte}\tsk}
\stick{\i{Mon Dieu!}\@ said he, embracing me\tsh}
\i{Vous etes Yorick}.

The Count instantly put the Shake-\break
\stick{spear into his pocket\tsk and left me}
alone in his room.

\patch[2\baselineskip]{\lsss{THE}}

\head{34pt}{68}{THE\enspace PASSPORT}
\head{22pt}{44}{VERSAILLES}
\vskip 4pt

\initial{I}{\ls{COULD}} not conceive why the\break
Count de B\anon\ had gone so\break
\stick{abruptly out of the room, any more}
\stick{than I could conceive why he had}
\stick{put the Shakespear into his pocket\tskk}
\stick{\i{Mysteries which must explain themselves},}
\stick{\i{are not worth the loss of time, which a}}
\stck{\i{conjecture about them takes up:} ’twas bet-}
\stick{ter to read Shakespear; so taking up}
\stick{\lqq \i{Much Ado about Nothing},\rqq\ I trans-}
\stick{ported myself instantly from the chair}
\stick{I sat in to Messina in Sicily, and}
\stick{got so busy with Don Pedro and Be-}
\catch{nedick}
\stick{nedick and Beatrice, that I thought}
\stick{not of Versailles, the Count, or the}
Passport.

Sweet pliability of man’s spirit,\break
\stick{that can at once surrender itself to}
\stick{illusions, which cheat expectation and}
\stick{sorrow of their weary moments!\tskk}
\stick{\tskk long\tskk long since had ye number’d}
\stick{out my days, had I not trod so great}
\stick{a part of them upon this enchanted}
\stick{ground: when my way is too rough}
\stick{for my feet, or too steep for my}
\stick{strength, I get off it, to some smooth}
\stick{velvet path which fancy has scattered}
\stick{over with rose-buds of delights; and}
\stick{having taken a few turns in it, come}
\stick{back strengthen’d and refresh’d\tskk}
\stick{When evils press sore upon me, and}
\catch{there}
\stick{there is no retreat from them in this}
\stick{world, then I take a new course\tskk I}
\stick{leave it\tskk and as I have a clearer idea of}
\stick{the elysian fields than I have of heaven,}
\stick{I force myself, like Eneas, into them}
\stick{\tskk I see him meet the pensive shade of}
\stick{his forsaken Dido\tskk and wish to re-}
\stick{cognize it\tskk I see the injured spirit}
\stick{wave her head, and turn off silent}
\stick{from the author of her miseries and}
\stick{dishonours\tskk I lose the feelings for}
\stick{myself in hers\tskk and in those af-}
\stick{fections which were wont to make}
\stick{me mourn for her when I was at}
school.

\i{Surely this is not walking in a vain}\break
\stick{\i{shadow\tskk nor does man disquiet himself}}
\stick{in vain, \i{by it}\tskk he oftener does so in}
\catch{trusting}
\stick{trusting the issue of his commotions}
\stick{to reason only.\tskk I can safely say for}
\stick{myself, I was never able to conquer}
\stick{any one single bad sensation in my}
\stick{heart so decisively, as by beating}
\stick{up as fast as I could for some kindly}
\stick{and gentle sensation, to fight it upon}
its own ground.

When I had got to the end of the\break
\stick{third act, the Count de B\anon\ en-}
\stick{tered with my Passport in his hand.}
\stick{Mons.\@ le Duc de C\anon, said the}
\stick{Count, is as good a prophet, I dare}
\stick{say, as he is a statesman\tskk \i{Un homme}}
\stick{\i{qui rit}, said the duke, \i{ne sera jamais}}
\stick{\i{dangereuz}.\tskk Had it been for any one}
\stick{but the king’s jester, added the Count, I}
\catch{could}
\stick{could not have got it these two}
\stick{hours.\tsk \i{Pardonnez moi}, Mons.\@ Le}
\stick{Compte, said I\tskk I am not the king’s}
\stick{jester.\tskk But you are Yorick?\tskk Yes.\tskk}
\stick{\i{Et vous plaisantez?}\tskk I answered, In-}
\stick{deed I did jest\tskk but was not paid for}
\stick{it\tskk ’twas entirely at my own ex-}
pence.

We have no jester at court, Mons.\break
\stick{Le Compte, said I, the last we had}
\stick{was in the licentious reign of Charles}
\stck{the IId\tskk since which time our manners}
\stick{have been so gradually refining,}
\stick{that our court at present is so full of}
\stick{patriots, who wish for \i{nothing} but}
\stick{the honours and wealth of their}
\stick{country\tskk and our ladies are all so}
\stick{chaste, so spotless, so good, so devout}
\catch{\tskk there}
\stick{\tskk there is nothing for a jester to make}
a jest of\tskk

\i{Voila un persiflage!}\@ cried the\break
Count.


\patch[13\baselineskip]{\lsss{THE}}
\head{34pt}{68}{THE\enspace PASSPORT}
\head{22pt}{44}{VERSAILLES}
\vskip 2pt

\dropcap{A}{s} the Passport was directed to\break 
all lieutenant governors, go-\break
\stick{vernors, and commandants of cities,}
\stick{generals of armies, justiciaries, and}
\stick{all officers of justice, to let Mr.\@ Yo-}
\stick{rick, the king’s jester, and his baggage,}
\stick{travel quietly along\tskk I own the tri-}
\stick{umph of obtaining the Passport was}
\stick{not a little tarnish’d by the figure I}
\stick{cut in it.\tskk But there is nothing un-}
\stick{mixt in this world; and some of the}
\stick{gravest of our divines have carried}
\stick{it so far as to affirm, that enjoyment it-}
\catch{self}
\stick{self was attended even with a sigh\tskk}
\stick{and that the greatest \i{they knew of},}
\stick{terminated \i{in a general way}, in little}
better than a convulsion.

I remember the grave and learned\break
\stick{Bevoriskius, in his commentary upon}
\stick{the generations from Adam, very}
\stick{naturally breaks off in the middle of}
\stick{a note to give an account to the}
\stick{world of a couple of sparrows upon}
\stick{the out-edge of his window, which}
\stick{had incommoded him all the time he}
\stick{wrote, and at last had entirely taken}
him off from his genealogy.

\tskk ’Tis strange! writes Bevoriskius;\break
\stick{but the facts are certain, for I have}
\catch{had}
\stick{had the curiosity to mark them}
\stick{down one by one with my pen\tskk but}
\stick{the cocks-sparrow, during the little}
\stick{time that I could have finished the}
\stick{other half this note, has actually}
\stick{interrupted me with the reiteration}
\stick{of his caresses three and twenty times}
and a half.

How merciful, adds Bevoriskius,\break
is heaven to his creatures!

Ill fated Yorick! that the gravest\break
\stick{of thy brethren should be able to}
\stick{write that to the world, which stains}
\stick{thy face with crimson, to copy in}
even thy study.

\patchv[\baselineskip]{G}{But}

But this is nothing to my travels\break
\stick{\tskk So I twice\tskk twice beg pardon}
for it.

\patch[15\baselineskip]{\lss{CHA}-}



\head{32pt}{76}{CHARACTER}
\head{22pt}{44}{VERSAILLES}

\vskip 2pt

\dropcap[findent=-0.5pt]{A}{nd} how do you find the\break
French? said the Count de\break
\stick{B\anon, after he had given me the}
Passport.

The reader may suppose that after\break
\stick{so obliging a proof of courtesy, I}
\stick{could not be at a loss to say some-}
thing handsome to the enquiry.

\tskk \i{Mais passe, pour cela}\tskk Speak\break
\stick{frankly, said he; do you find all the}
\stick{urbanity in the French which the}
\catchs{G 2}{world}
\stick{world give us the honour of?\tskk I had}
\stick{found every thing, I said, which con-}
\stick{firmed it\tskk \i{Vraiment}, said the count.}
\stick{\tskk\i{Les Francois sont poli.s}\sic\tskk To an ex-}
cess, replied I.

The count took notice of the word\break
\stick{\i{excesse}; and would have it I meant}
\stick{more than I said.  I defended myself}
\stick{a long time as well as I could}
\stick{against it\tskk he insisted I had a reserve,}
\stick{and that I would speak my opinion}
frankly.

I believe, Mons.\@ Le Compte, said\break
\stick{I, that man has a certain compass, as}
\stick{well as an instrument; and that the}
\stick{social and other calls have occasion}
\stick{by turns for every key in him; so}
\catch{that}
\stick{that if you begin a note too high}
\stick{or too low, there must be a want}
\stick{either in the upper or under part,}
\stick{to fill up the system of harmony.\tskk}
\stick{The Count de B\anon\ did not un-}
\stick{derstand music, so desired me to ex-}
\stick{plain it some other way. A polish’d}
\stick{nation, my dear Count, said I, makes}
\stick{every one its debtors; and besides}
\stick{urbanity itself, like the fair sex, has}
\stick{so many charms; it goes against the\sic}
\stick{the heart to say it can do ill; and}
\stick{yet, I believe, there is but a certain}
\stick{line of perfection, that man, take him}
\stick{altogether, is empower’d to arrive at}
\stick{\tskk if he gets beyond, he rather ex-}
\stick{changes qualities, than gets them.}
\stick{I must not presume to say, how far this}
\stick{has affected the French in the sub-}
\catchs{G 3}{ject}
\stick{ject we are speaking of\tskk but, should}
\stick{it ever be the case of the English, in}
\stick{the progress of their refinements, to}
\stick{arrive at the same polish which dis-}
\stick{tinguishes the French, if we did not}
\stick{lose the \i{politesse du cœur}, which in-}
\stick{clines men more to human actions,}
\stick{than courteous ones\tskk we should at}
\stick{least lose that distinct variety and}
\stick{originality of character, which dis-}
\stick{tinguishes them, not only from each}
\stick{other, but from all the world be-}
sides.

I had a few king William’s shil-\break
\stick{lings as smooth as glass in my pocket;}
\stick{and foreseeing they would be of use}
\stick{in the illustration of my hypothesis,}
\catch{I had}
\stick{I had got them into my hand, when}
I had proceeded so far\tskk

See, Mons.\@ Le Compte, said I,\break
\stick{rising up, and laying them before}
\stick{him upon the table\tskk by jingling and}
\stick{ribbing one against another for seventy}
\stick{years together in one body’s pocket}
\stick{or another’s, they are become so}
\stick{much alike, you can scarce distin-}
guish one shilling from another.

The English, like antient medals,\break
\stick{kept more apart, and passing but few}
\stick{peoples hands, preserve the first}
\stick{sharpnesses which the fine hand of na-}
\stick{ture has given them\tskk they are not}
\stick{so pleasant to feel\tskk but in return,}
\catchs{G 4}{the}
\stick{the legend is so visible, that at the}
\stick{first look you see whose image and}
\stick{superscription they bear.\tskk But the}
\stick{French, Mons.\@ Le Compte, added I,}
\stick{wishing to soften what I had said, have}
\stick{so many excellences, they can the}
\stick{better spare this\tskk they are a loyal, a}
\stick{gallant, a generous, an ingenious,}
\stick{and good temper’d people as is}
\stick{under heaven\tskk if they have a fault}
\tskk they are too \i{serious}.

\i{Mon Dieu!} cried the Count, rising\break
out of his chair.

\i{Mais vous plaisantez}, said he, cor-\break
\stick{recting his exclamation.\tskk I laid my}
\stick{hand upon my breast, and with ear-}
\catch{nest}
\stick{nest gravity assured him it was my}
most settled opinion.

The Count said he was mortified,\break
\stick{he could not stay to hear my rea-}
\stick{sons, being engaged to go that mo-}
\stick{ment to dine with the Duc de}
C\anon.

But if it is not too far to come to\break
\stick{Versailles to eat your soup with me,}
\stick{I beg, before you leave France, I may}
\stick{have the pleasure of knowing you}
\stick{retract your opinion\tskk or, in what}
\stick{manner you support it.\tskk But if you}
\stick{do support it, Mons.\@ Anglois, said}
\stick{he, you must do it with all your}
\stick{powers, because you have the whole}
\catch{world}
\stick{world against you.\tskk I promised the}
\stick{Count I would do myself the honour}
\stick{of dining with him before I set out}
for Italy\tskk so took my leave.

\patch[13\baselineskip]{\lsss{THE}}


\head{24pt}{48}{THE\enspace TEMPTATION}
\head{24pt}{64}{PARIS}
\vskip 0pt

\dropcap{W}{hen} I alighted at the hotel,\break
the porter told me a young\break
\stick{woman with a band-box had been that}
\stick{moment enquiring for me.\tskk I do not}
\stick{know, said the porter, whether she}
\stick{is gone away or no.  I took the key}
\stick{of my chamber of him, and went}
\stick{up stairs; and when I had got with-}
\stick{in ten steps of the top of the landing}
\stick{before my door, I met her coming}
easily down.

It was the fair \i{fille de chambre} I\break
\stick{had walked along the Quai de Conti}
\catch{with:}
\stick{with: Madame de R\anon\ had sent}
\stick{her upon some commissions to a \i{mer}-}
\stick{\i{chande des modes} within a step or two}
\stick{of the hotel de Modene; and as I had}
\stick{fail’d in waiting upon her, had bid}
\stick{her enquire if I had left Paris; and}
\stick{if so, whether I had not left a letter}
address’d to her.

As the fair \i{fille de chambre} was so\break
\stick{near my door she returned back, and}
\stick{went into the room with me for a}
\stick{moment or two whilst I wrote a}
card.

It was a fine still evening in the\break
\stick{latter end of the month of May\tskk}
\stick{the crimson window curtains (which}
\stick{were of the same colour as those of}
\catch{the}
\stick{the bed) were drawn close\tskk the sun}
\stick{was setting, and reflected through}
\stick{them so warm a tint into the fair \i{fille}}
\stick{\i{de chambre}’s face\tskk I thought she}
\stick{blush’d\tskk the idea of it made me}
\stick{blush myself\tskk we were quite alone;}
\stick{and that super-induced a second blush}
before the first could get off.

There is a sort of a pleasing half\break
\stick{guilty blush, where the blood is}
\stick{more in fault than the man\tskk ’tis sent}
\stick{impetuous from the heart, and virtue}
\stick{flies after it\tskk not to call it back, but}
\stick{to make the sensation of it more}
\stick{delicious to the nerves\tskk ’tis asso-}
ciated.\tskk\sic

\patch[\baselineskip]{But}

But I’ll not describe it.\tskk I felt\break
\stick{something at first within me which}
\stick{was not in strict unison with the}
\stick{lesson of virtue I had given her}
\stick{the night before\tskk I sought five mi-}
\stick{nutes for a card\tskk I knew I had not}
\stick{one.\tskk I took up a pen\tskk I laid it}
\stick{down again\tskk my hand trembled\tskk}
the devil was in me.

I know as well as any one, he is an\break
\stick{adversary, whom if we resist, he will}
\stick{fly from us\tskk but I seldom resist him}
\stick{at all; from a terror, that though I}
\stick{may conquer, I may still get a hurt}
\stick{in the combat\tskk so I give up the}
\stick{triumph, for security; and instead}
\stick{of thinking to make him fly, I ge-}
nerally fly myself.

\patch{The}

The fair \i{fille de chambre} came close\break
\stick{up to the bureau where I was looking}
\stick{for a card\tskk took up first the pen I}
\stick{cast down, then offered to hold me}
\stick{the ink; she offer’d it so sweetly, I}
\stick{was going to accept it\tskk but I durst}
\stick{not\tskk I have nothing, my dear, said}
\stick{I, to write upon.\tskk Write it, said she,}
simply, upon any thing.\tskk

I was just going to cry out, Then\break
\stick{I will write it, fair girl!\@ upon thy}
lips.\tskk

If I do, said I, I shall perish\tskk\break
\stick{so I took her by the hand, and led}
\stick{her to the door, and begg’d she would}
\stick{not forget the lesson I had given her}
\stick{\tskk She said, Indeed she would not\tskk}
\catch{and}
\stick{and as she utter’d it with some ear-}
\stick{nestness, she turned about, and gave}
\stick{me both her hands, closed together,}
\stick{into mine\tskk it was impossible not to}
\stick{compress them in that situation\tskk I}
\stick{wish’d to let them go; and all the}
\stick{time I held them, I kept arguing}
\stick{within myself against it\tskk and still I}
\stick{held them on.\tskk In two minutes I}
\stick{found I had all the battle to fight}
\stick{over again\tskk and I felt my legs and}
\stick{every limb about me tremble at the}
idea.

The foot of the bed was within a\break
\stick{yard and a half of the place where}
\stick{we were standing\tskk I had still hold of}
\stick{her hands\tskk and how it happened I}
\stick{can give no account; but I neither}
\catch{ask’d}
\stick{ask’d her\tskk nor drew her\tskk nor did I}
\stick{think of the bed\tskk but so it did hap-}
pen, we both sat down.

I’ll just shew you, said the fair\break
\stick{\i{fille de chambre}, the little purse I have}
\stick{been making to-day to hold your}
\stick{crown.  So she put her hand into}
\stick{her right pocket, which was next me,}
\stick{and felt for it for some time\tskk then into}
\stick{into\sic the left\tskk \lqq She had lost it.\rqq \tskk I}
\stick{never bore expectation more quietly}
\stick{\tskk it was in her right pocket at last\tskk}
\stick{she pulled it out; it was of green}
\stick{taffeta, lined with a little bit of white}
\stick{quilted sattin, and just big enough}
\stick{to hold the crown\tskk she put it into}
\stick{my hand\tskk it was pretty; and I held it}
\stick{ten minutes with the back of my hand}
\catchv{H}{resting}
\stick{resting upon her lap\tskk looking some-}
\stick{times at the purse, sometimes on one}
side of it.

A stitch or two had broke out in\break
\stick{the gathers of my stock\tskk the fair}
\stick{\i{fille de chambre}, without saying a word,}
\stick{took out her little hussive, threaded}
\stick{a small needle, and sew’d it up\tskk}
\stick{I foresaw it would hazard the glory}
\stick{of the day; and as she passed her hand}
\stick{in silence across and across my neck}
\stick{in the manœuvre, I felt the laurels}
\stick{shake which fancy had wreath’d}
about my head.

A strap had given way in her\break
\stick{walk, and the buckle of her shoe was}
\stick{just falling off\tskk See, said the \i{fille de}}
\catch{\i{chambre},}
\stick{\i{chambre}, holding up her foot\tskk I}
\stick{could not for my soul but fasten the}
\stick{buckle in return, and putting in the}
\stick{strap\tskk and lifting up the other foot}
\stick{with it, when I had done, to see both}
\stick{were right\tskk in doing it too suddenly}
\stick{\tskk it unavoidably threw the fair \i{fille de}}
\i{chambre} off her center\tskk and then\tskk

\patchs[10\baselineskip]{H 2}{\lsss{THE}}


\head{32pt}{48}{THE CONQUEST}

\vskip 6pt


\dropcap{Y}{es}\tshh and then\tskk Ye whose\break
clay-cold heads and luke-warm\break
\stick{hearts can argue down or mask your}
\stick{passions\tskk tell me, what trespass is it}
\stick{that man should have them? or how}
\stick{his spirit stands answerable, to the}
\stick{father of spirits, but for his conduct}
under them?

If nature has so wove her web of\break
\stick{kindness, that some threads of love}
\stick{and desire are entangled with the}
\stick{piece\tskk must the whole web be rent}
\stick{in drawing them out?\tskk Whip me such}
\stick{stoics, great governor of nature!}
\catch{said}
\stick{said I to myself\tskk Wherever thy pro-}
\stick{vidence shall place me for the trials of}
\stick{my virtue\tskk whatever is my danger}
\stick{\tskk whatever is my situation\tskk let me}
\stick{feel the movements which rise out of}
\stick{it, and which belong to me as a man}
\stick{\tskk and if I govern them as a good}
\stick{one\tskk I will trust the issues to thy}
\stick{justice, for thou hast made us\tskk and}
not we ourselves.

As I finish’d my address, I raised\break
\stick{the fair \i{fille de chambre} up by the}
\stick{hand, and led her out of the room}
\stick{\tskk she stood by me till I lock’d the}
\stick{door and put the key in my pocket}
\stick{\tskk \i{and then}\tskk the victory being}
\stick{quite decisive\tskk and not till then, I}
\catchs{H 3}{press’d}
\stick{press’d my lips to her cheek, and,}
\stick{taking her by the hand again,}
\stick{led her safe to the gate of the}
hotel.

\patch[14\baselineskip]{\lsss{THE}}

\head{32pt}{48}{THE MYSTERY}
\head{24pt}{56}{PARIS}

\vskip 6pt


\dropcap{I}{f} a man knows the heart, he\break
will know it was impossible to go\break
\stick{back instantly to my chamber\tskk it}
\stick{was touching a cold key with a flat}
\stick{third to it, upon the close of a piece}
\stick{of musick, which had call’d forth my}
\stick{affections\tskk therefore, when I let go}
\stick{the hand of the \i{fille de chambre}, I}
\stick{remain’d at the gate of the hotel for}
\stick{some time, looking at every one who}
\stick{pass’d by, and forming conjectures}
\stick{upon them, till my attention got fix’d}
\catchs{H 4}{upon}
\stick{upon a single object which confounded}
all kind of reasoning upon him.

It was a tall figure of a philoso-\break
\stick{phic serious, adust look, which}
\stick{pass’d and repass’d sedately along the}
\stick{street, making a turn of about sixty}
\stick{paces on each side of the gate of the}
\stick{hotel\tskk the man was about fifty-two}
\stick{\tskk had a small cane under his arm\tskk}
\stick{was dress’d in a dark drab-colour’d}
\stick{coat, waistcoat, and breeches, which}
\stick{seem’d to have seen some years ser-}
\stick{vice\tskk they were still clean, and there}
\stick{was a little air of frugal \i{propretè}}
\stick{throughout him.  By his pulling off}
\stick{his hat, and his attitude of accosting}
\stick{a good many in his way, I saw he}
\stick{was asking charity; so I got a sous or}
\catch{two}
\stick{two out of my pocket ready to give}
\stick{him, as he took me in his turn\tskk he}
\stick{pass’d by me without asking any}
\stick{thing\tskk and yet did not go five steps}
\stick{further before he ask’d charity of a}
\stick{little woman\tskk I was much more}
\stick{likely to have given of the two\tskk He}
\stick{had scarce done with the woman,}
\stick{when he pull’d off his hat to another}
\stick{who was coming the same way.\tskk An}
\stick{ancient gentleman came slowly\tskk and,}
\stick{after him, a young smart one\tskk He}
\stick{let them both pass, and ask’d nothing:}
\stick{I stood observing him half an hour, in}
\stick{which time he had made a dozen}
\stick{turns backwards and forwards, and}
\stick{found that he invariably pursued the}
same plan.

\patch[\baselineskip]{There}

There were two things very sin-\break
\stick{gular in this, which set my brain to}
\stick{work, and to no purpose\tskk the first}
\stick{was, why the man should \i{only} tell}
\stick{his story to the sex\tskk and secondly\tskk}
\stick{what kind of story it was, and what}
\stick{species of eloquence it could be,}
\stick{which soften’d the hearts of the}
\stick{women, which he knew ’twas to no}
purpose to practise upon the men.

There were two other circum-\break
\stick{stances which entangled this mystery}
\stick{\tskk the one was, he told every woman}
\stick{what he had to say in her ear, and in}
\stick{a way which had much more the air}
\stick{of a secret than a petition\tskk the}
\stick{other was, it was always successful}
\stick{\tskk he never stopp’d a woman, but she}
\catch{pull’d}
\stick{pull’d out her purse, and immedi-}
ately gave him something.

I could form no system to explain\break
the phenomenon.

I had got a riddle to amuse me for\break
\stick{the rest of the evening, so I walk’d}
up stairs to my chamber.


\patch[9\baselineskip]{\lsss{\scshape The}}



\head{32pt}{48}{THE CASE OF CONSCIENCE}
\head{24pt}{48}{PARIS}

\vskip 6pt


\initial{I}{was} immediately followed up
by the master of the hotel, who came into
my room to tell me I must provide lodgings
elsewhere.\tskk How so, friend?  said
I.\tskk He answered, I had had a young
woman lock’d up with me two hours that
evening in my bedchamber, and ’twas
against the rules of his house.\tskk Very
well, said I, we’ll all part friends
then\tskk for the girl is no worse,\tskk
and I am no worse\tskk and you will be
just as I found you.\tskk It was enough,
he said, to overthrow the credit of his
hotel.\tskk \i{Voyez vous}, Monsieur, said
he, pointing to the foot of the bed we had
been sitting upon.\tskk I own it had
something of the appearance of an
evidence; but my pride not suffering me to
enter into any detail of the case, I
exhorted him to let his soul sleep in
peace, as I resolved to let mine do that
night, and that I would discharge what I
owed him at breakfast.

I should not have minded, Monsieur, said
he, if you had had twenty girls\tskk ’Tis
a score more, replied I, interrupting him,
than I ever reckon’d upon\tskk Provided,
added he, it had been but in a
morning.\tskk And does the difference of
the time of the day at Paris make a
difference in the sin?\tskk It made a
difference, he said, in the scandal.\tskk
I like a good distinction in my heart; and
cannot say I was intolerably out of temper
with the man.\tskk I own it is necessary,
resumed the master of the hotel, that a
stranger at Paris should have the
opportunities presented to him of buying
lace and silk stockings and ruffles, \i{et
tout cela}\tskk and ’tis nothing if a
woman comes with a band-box.\tskk O, my
conscience! said I, she had one but I
never look’d into it.\tskk Then Monsieur,
said he, has bought nothing?\tskk Not one
earthly thing, replied I.\tskk Because,
said he, I could recommend one to you who
would use you \i{en conscience}.\tskk But
I must see her this night, said I.\tskk He
made me a low bow, and walk’d down.

Now shall I triumph over this \i{maître
d’hôtel}, cried I\tskk and what then?
Then I shall let him see I know he is a
dirty fellow.\tskk And what then?  What
then?\tskk I was too near myself to say it
was for the sake of others.\tskk I had no
good answer left\tskk there was more of
spleen than principle in my project, and I
was sick of it before the execution.

In a few minutes the grisette came in with
her box of lace.\tskk I’ll buy nothing,
however, said I, within myself.

The grisette would shew me
everything.\tskk I was hard to please: she
would not seem to see it; she opened her
little magazine, and laid all her laces
one after another before me\tskk unfolded
and folded them up again one by one with
the most patient sweetness.\tskk I might
buy\tskk or not;\tskk she would let me
have everything at my own price\tskk the
poor creature seem’d anxious to get a
penny; and laid herself out to win me, and
not so much in a manner which seem’d
artful, as in one I felt simple and
caressing.

If there is not a fund of honest
gullibility in man, so much the
worse\tskk my heart relented, and I gave
up my second resolution as quietly as the
first.\tskk Why should I chastise one for
the trespass of another?  If thou art
tributary to this tyrant of an host,
thought I, looking up in her face, so much
harder is thy bread.

If I had not had more than four louis
d’ors in my purse, there was no such thing
as rising up and shewing her the door,
till I had first laid three of them out in
a pair of ruffles.

\tskk The master of the hotel will share
the profit with her\tskk no matter,\tskk
then I have only paid as many a poor soul
has \i{paid} before me, for an act he
\i{could} not do, or think of.




\head{32pt}{48}{THE RIDDLE}
\head{24pt}{48}{PARIS}

\vskip 6pt


\dropcap{W}{hen} La Fleur came up to wait
upon me at supper, he told me how sorry
the master of the hotel was for his
affront to me in bidding me change my
lodgings.

A man who values a good night’s rest will
not lie down with enmity in his heart, if
he can help it.\tskk So I bid La Fleur
tell the master of the hotel, that I was
sorry on my side for the occasion I had
given him\tskk and you may tell him, if
you will, La Fleur, added I, that if the
young woman should call again, I shall not
see her.

This was a sacrifice not to him, but
myself, having resolved, after so narrow
an escape, to run no more risks, but to
leave Paris, if it was possible, with all
the virtue I enter’d it.

\i{C’est déroger à noblesse},
\i{Monsieur}, said La Fleur, making me a
bow down to the ground as he said it.\tskk
\i{Et encore}, \i{Monsieur}, said he, may
change his sentiments\tskk and if (\i{par
hazard}) he should like to amuse
himself\tskk I find no amusement in it,
said I, interrupting him.\tskk 

\i{Mon Dieu}! said La Fleur\tskk and took
away.

In an hour’s time he came to put me to
bed, and was more than commonly
officious\tskk something hung upon his
lips to say to me, or ask me, which he
could not get off: I could not conceive
what it was, and indeed gave myself little
trouble to find it out, as I had another
riddle so much more interesting upon my
mind, which was that of the man’s asking
charity before the door of the hotel.\tskk
I would have given anything to have got to
the bottom of it; and that, not out of
curiosity\tskk ’tis so low a principle of
enquiry, in general, I would not purchase
the gratification of it with a two-sous
piece\tskk but a secret, I thought, which
so soon and so certainly soften’d the
heart of every woman you came near, was a
secret at least equal to the philosopher’s
stone; had I both the Indies, I would have
given up one to have been master of it.

I toss’d and turn’d it almost all night
long in my brains to no manner of purpose;
and when I awoke in the morning, I found
my spirits as much troubled with my
dreams, as ever the King of Babylon had
been with his; and I will not hesitate to
affirm, it would have puzzled all the wise
men of Paris as much as those of Chaldea
to have given its interpretation.




\head{32pt}{48}{LE DIMANCHE}
\head{24pt}{48}{PARIS}

\vskip 6pt


\dropcap{I}{t} was Sunday; and when La
Fleur came in, in the morning, with my
coffee and roll and butter, he had got
himself so gallantly array’d, I scarce
knew him.

I had covenanted at Montreuil to give him
a new hat with a silver button and loop,
and four louis d’ors, \i{pour s’adoniser},
when we got to Paris; and the poor fellow,
to do him justice, had done wonders with
it.

He had bought a bright, clean, good
scarlet coat, and a pair of breeches of
the same.\tskk They were not a crown
worse, he said, for the wearing.\tskk I
wish’d him hang’d for telling me.\tskk
They look’d so fresh, that though I knew
the thing could not be done, yet I would
rather have imposed upon my fancy with
thinking I had bought them new for the
fellow, than that they had come out of the
Rue de Friperie.

This is a nicety which makes not the heart
sore at Paris.

He had purchased, moreover, a handsome
blue satin waistcoat, fancifully enough
embroidered\tskk this was indeed
something the worse for the service it had
done, but ’twas clean scour’d\tskk the
gold had been touch’d up, and upon the
whole was rather shewy than
otherwise\tskk and as the blue was not
violent, it suited with the coat and
breeches very well: he had squeez’d out of
the money, moreover, a new bag and a
solitaire; and had insisted with the
\i{fripier} upon a gold pair of garters to
his breeches knees.\tskk He had purchased
muslin ruffles, \i{bien brodées}, with
four livres of his own money\tskk and a
pair of white silk stockings for five
more\tskk and to top all, nature had
given him a handsome figure, without
costing him a sous.

He entered the room thus set off, with his
hair dressed in the first style, and with
a handsome bouquet in his breast.\tskk In
a word, there was that look of festivity
in everything about him, which at once put
me in mind it was Sunday\tskk and, by
combining both together, it instantly
struck me, that the favour he wish’d to
ask of me the night before, was to spend
the day as every body in Paris spent it
besides.  I had scarce made the
conjecture, when La Fleur, with infinite
humility, but with a look of trust, as if
I should not refuse him, begg’d I would
grant him the day, \i{pour faire le galant
vis-à-vis de sa maîtresse}.

Now it was the very thing I intended to do
myself vis-à-vis Madame de R\anon .\tskk I
had retained the remise on purpose for it,
and it would not have mortified my vanity
to have had a servant so well dress’d as
La Fleur was, to have got up behind it: I
never could have worse spared him.

But we must \i{feel}, not argue in these
embarrassments.\tskk The sons and
daughters of Service part with liberty,
but not with nature, in their contracts;
they are flesh and blood, and have their
little vanities and wishes in the midst of
the house of bondage, as well as their
task-masters\tskk no doubt, they have set
their self-denials at a price\tskk and
their expectations are so unreasonable,
that I would often disappoint them, but
that their condition puts it so much in my
power to do it.

\i{Behold}\tskk \i{Behold}, \i{I am thy
servant}\tskk disarms me at once of the
powers of a master.\tskk 

Thou shalt go, La Fleur! said I.

\tskk And what mistress, La Fleur, said I,
canst thou have picked up in so little a
time at Paris?  La Fleur laid his hand
upon his breast, and said ’twas a
\i{petite demoiselle}, at Monsieur le
Count de B\anon’s.\tskk La Fleur had a
heart made for society; and, to speak the
truth of him, let as few occasions slip
him as his master\tskk so that somehow or
other\tskk but how,\tskk heaven
knows\tskk he had connected himself with
the demoiselle upon the landing of the
staircase, during the time I was taken up
with my passport; and as there was time
enough for me to win the Count to my
interest, La Fleur had contrived to make
it do to win the maid to his.  The family,
it seems, was to be at Paris that day, and
he had made a party with her, and two or
three more of the Count’s household, upon
the boulevards.

Happy people! that once a week at least
are sure to lay down all your cares
together, and dance and sing and sport
away the weights of grievance, which bow
down the spirit of other nations to the
earth.




\head{32pt}{48}{THE FRAGMENT}
\head{24pt}{48}{PARIS}

\vskip 6pt


\dropcap{L}{a Fleur} had left me something
to amuse myself with for the day more than
I had bargain’d for, or could have enter’d
either into his head or mine.

He had brought the little print of butter
upon a currant leaf: and as the morning
was warm, and he had a good step to bring
it, he had begg’d a sheet of waste paper
to put betwixt the currant leaf and his
hand.\tskk As that was plate sufficient, I
bade him lay it upon the table as it was;
and as I resolved to stay within all day,
I ordered him to call upon the
\i{traîteur}, to bespeak my dinner, and
leave me to breakfast by myself.

When I had finished the butter, I threw
the currant-leaf out of the window, and
was going to do the same by the waste
paper\tskk but stopping to read a line
first, and that drawing me on to a second
and third\tskk I thought it better worth;
so I shut the window, and drawing a chair
up to it, I sat down to read it.

It was in the old French of Rabelais’s
time, and for aught I know might have been
wrote by him\tskk it was moreover in a
Gothic letter, and that so faded and gone
off by damps and length of time, it cost
me infinite trouble to make anything of
it.\tskk I threw it down; and then wrote a
letter to Eugenius\tskk then I took it up
again, and embroiled my patience with it
afresh\tskk and then to cure that, I
wrote a letter to Eliza.\tskk Still it
kept hold of me; and the difficulty of
understanding it increased but the desire.

I got my dinner; and after I had
enlightened my mind with a bottle of
Burgundy; I at it again\tskk and, after
two or three hours poring upon it, with
almost as deep attention as ever Gruter or
Jacob Spon did upon a nonsensical
inscription, I thought I made sense of it;
but to make sure of it, the best way, I
imagined, was to turn it into English, and
see how it would look then\tskk so I went
on leisurely, as a trifling man does,
sometimes writing a sentence\tskk then
taking a turn or two\tskk and then
looking how the world went, out of the
window; so that it was nine o’clock at
night before I had done it.\tskk I then
began and read it as follows.




\head{32pt}{48}{THE FRAGMENT}
\head{24pt}{48}{PARIS}

\vskip 6pt


\tskk NOW, as the notary’s wife disputed
the point with the notary with too much
heat\tskk I wish, said the notary,
(throwing down the parchment) that there
was another notary here only to set down
and attest all this.\tskk 

\tskk And what would you do then,
Monsieur? said she, rising hastily
up.\tskk The notary’s wife was a little
fume of a woman, and the notary thought it
well to avoid a hurricane by a mild
reply.\tskk I would go, answered he, to
bed.\tskk You may go to the devil,
answer’d the notary’s wife.

Now there happening to be but one bed in
the house, the other two rooms being
unfurnished, as is the custom at Paris,
and the notary not caring to lie in the
same bed with a woman who had but that
moment sent him pell mell to the devil,
went forth with his hat and cane and short
cloak, the night being very windy, and
walk’d out, ill at ease, towards the Pont
Neuf.

Of all the bridges which ever were built,
the whole world who have pass’d over the
Pont Neuf must own, that it is the
noblest\tskk the finest,\tskk the
grandest\tskk the lightest,\tskk the
longest\tskk the broadest, that ever
conjoin’d land and land together upon the
face of the terraqueous globe.

[\i{By this it seems as if the author of
the fragment had not been a Frenchman}.]

The worst fault which divines and the
doctors of the Sorbonne can allege against
it is, that if there is but a capfull of
wind in or about Paris, ’tis more
blasphemously \i{sacre Dieu’d} there than
in any other aperture of the whole
city\tskk and with reason good and
cogent, Messieurs; for it comes against
you without crying \i{garde d’eau}, and
with such unpremeditable puffs, that of
the few who cross it with their hats on,
not one in fifty but hazards two livres
and a half, which is its full worth.

The poor notary, just as he was passing by
the sentry, instinctively clapp’d his cane
to the side of it, but in raising it up,
the point of his cane catching hold of the
loop of the sentinel’s hat, hoisted it
over the spikes of the ballustrade clear
into the Seine.\tskk 

\tskk ’\i{Tis an ill wind}, said a
boatman, who catched it, \i{which blows
nobody any good}.

The sentry, being a Gascon, incontinently
twirled up his whiskers, and levell’d his
arquebuss.

Arquebusses in those days went off with
matches; and an old woman’s paper lantern
at the end of the bridge happening to be
blown out, she had borrow’d the sentry’s
match to light it\tskk it gave a moment’s
time for the Gascon’s blood to run cool,
and turn the accident better to his
advantage.\tskk ’\i{Tis an ill wind}, said
he, catching off the notary’s castor, and
legitimating the capture with the
boatman’s adage.

The poor notary crossed the bridge, and
passing along the Rue de Dauphine into the
fauxbourgs of St. Germain, lamented
himself as he walked along in this
manner\tskk 

Luckless man that I am! said the notary,
to be the sport of hurricanes all my
days\tskk to be born to have the storm of
ill language levell’d against me and my
profession wherever I go; to be forced
into marriage by the thunder of the church
to a tempest of a woman\tskk to be driven
forth out of my house by domestic winds,
and despoil’d of my castor by pontific
ones!\tskk to be here, bareheaded, in a
windy night, at the mercy of the ebbs and
flows of accidents!\tskk Where am I to lay
my head?\tskk Miserable man! what wind in
the two-and-thirty points of the whole
compass can blow unto thee, as it does to
the rest of thy fellow-creatures, good?

As the notary was passing on by a dark
passage, complaining in this sort, a voice
call’d out to a girl, to bid her run for
the next notary.\tskk Now the notary being
the next, and availing himself of his
situation, walk’d up the passage to the
door, and passing through an old sort of a
saloon, was usher’d into a large chamber,
dismantled of everything but a long
military pike\tskk a breastplate,\tskk a
rusty old sword, and bandoleer, hung up,
equidistant, in four different places
against the wall.

An old personage who had heretofore been a
gentleman, and unless decay of fortune
taints the blood along with it, was a
gentleman at that time, lay supporting his
head upon his hand in his bed; a little
table with a taper burning was set close
beside it, and close by the table was
placed a chair\tskk the notary sat him
down in it; and pulling out his inkhorn
and a sheet or two of paper which he had
in his pocket, he placed them before him;
and dipping his pen in his ink, and
leaning his breast over the table, he
disposed everything to make the
gentleman’s last will and testament.

Alas!  \i{Monsieur le Notaire}, said the
gentleman, raising himself up a little, I
have nothing to bequeath, which will pay
the expense of bequeathing, except the
history of myself, which I could not die
in peace, unless I left it as a legacy to
the world: the profits arising out of it I
bequeath to you for the pains of taking it
from me.\tskk It is a story so uncommon,
it must be read by all mankind\tskk it
will make the fortunes of your house.\tskk
The notary dipp’d his pen into his
inkhorn.\tskk Almighty Director of every
event in my life! said the old gentleman,
looking up earnestly, and raising his
hands towards heaven\tskk Thou, whose
hand has led me on through such a
labyrinth of strange passages down into
this scene of desolation, assist the
decaying memory of an old, infirm, and
broken-hearted man\tskk direct my tongue
by the spirit of thy eternal truth, that
this stranger may set down nought but what
is written in that BOOK, from whose
records, said he, clasping his hands
together, I am to be condemn’d or
acquitted!\tskk the notary held up the
point of his pen betwixt the taper and his
eye.\tskk 

It is a story, \i{Monsieur le Notaire},
said the gentleman, which will rouse up
every affection in nature\tskk it will
kill the humane, and touch the heart of
Cruelty herself with pity.\tskk 

\tskk The notary was inflamed with a
desire to begin, and put his pen a third
time into his ink-horn\tskk and the old
gentleman, turning a little more towards
the notary, began to dictate his story in
these words\tskk 

\tskk And where is the rest of it, La
Fleur? said I, as he just then enter’d the
room.




\head{32pt}{24}{THE FRAGMENT, AND THE
BOUQUET\fnast} \head{24pt}{48}{PARIS}

\dropcap{W}{hen} La Fleur came up close to
the table, and was made to comprehend what
I wanted, he told me there were only two
other sheets of it, which he had wrapped
round the stalks of a bouquet to keep it
together, which he had presented to the
demoiselle upon the boulevards.\tskk Then
prithee, La Fleur, said I, step back to
her to the Count de B\anon’s hotel, and
see if thou canst get it.\tskk There is no
doubt of it, said La Fleur\tskk and away
he flew.

In a very little time the poor fellow came
back quite out of breath, with deeper
marks of disappointment in his looks than
could arise from the simple irreparability
of the fragment.  \i{Juste Ciel}! in less
than two minutes that the poor fellow had
taken his last tender farewell of her\tskk
his faithless mistress had given his
\i{gage d’amour} to one of the Count’s
footmen\tskk the footman to a young
sempstress\tskk and the sempstress to a
fiddler, with my fragment at the end of
it.\tskk Our misfortunes were involved
together\tskk I gave a sigh,\tskk and La
Fleur echoed it back again to my ear.

\stick{\hfill\footnotesize\fnast Nosegay.\hfill}

\tskk How perfidious! cried La Fleur.\tskk
How unlucky! said I.

\tskk I should not have been mortified,
Monsieur, quoth La Fleur, if she had lost
it.\tskk Nor I, La Fleur, said I, had I
found it.

Whether I did or no will be seen
hereafter.




\head{32pt}{48}{THE ACT OF CHARITY}
\head{24pt}{48}{PARIS}

\vskip 6pt


\dropcap{T}{he} man who either disdains or
fears to walk up a dark entry may be an
excellent good man, and fit for a hundred
things, but he will not do to make a good
Sentimental Traveller.\tskk I count little
of the many things I see pass at broad
noonday, in large and open streets.\tskk
Nature is shy, and hates to act before
spectators; but in such an unobserved
corner you sometimes see a single short
scene of hers worth all the sentiments of
a dozen French plays compounded
together\tskk and yet they are absolutely
fine\tskk and whenever I have a more
brilliant affair upon my hands than
common, as they suit a preacher just as
well as a hero, I generally make my sermon
out of ’em\tskk and for the text,\tskk
\lqq Cappadocia, Pontus and Asia, Phrygia and
Pamphylia,\rqq \tskk is as good as any one in
the Bible.

There is a long dark passage issuing out
from the Opera Comique into a narrow
street; ’tis trod by a few who humbly wait
for a \i{fiacre}, {649} or wish to get off
quietly o’foot when the opera is done.  At
the end of it, towards the theatre, ’tis
lighted by a small candle, the light of
which is almost lost before you get
half-way down, but near the door\tskk ’tis
more for ornament than use: you see it as
a fixed star of the least magnitude; it
burns\tskk but does little good to the
world, that we know of.

{649}  Hackney coach.

In returning along this passage, I
discerned, as I approached within five or
six paces of the door, two ladies standing
arm-in-arm with their backs against the
wall, waiting, as I imagined, for a
\i{fiacre}\tskk as they were next the
door, I thought they had a prior right; so
edged myself up within a yard or little
more of them, and quietly took my
stand.\tskk I was in black, and scarce
seen.

The lady next me was a tall lean figure of
a woman, of about thirty-six; the other of
the same size and make, of about forty:
there was no mark of wife or widow in any
one part of either of them\tskk they
seem’d to be two upright vestal sisters,
unsapped by caresses, unbroke in upon by
tender salutations.\tskk I could have
wish’d to have made them happy\tskk their
happiness was destin’d that night, to come
from another quarter.

A low voice, with a good turn of
expression, and sweet cadence at the end
of it, begg’d for a twelve-sous piece
betwixt them, for the love of heaven.  I
thought it singular that a beggar should
fix the quota of an alms\tskk and that the
sum should be twelve times as much as what
is usually given in the dark.\tskk They
both seemed astonished at it as much as
myself.\tskk Twelve sous! said one.\tskk A
twelve-sous piece! said the other\tskk
and made no reply.

The poor man said, he knew not how to ask
less of ladies of their rank; and bow’d
down his head to the ground.

Poo! said they\tskk we have no money.

The beggar remained silent for a moment or
two, and renew’d his supplication.

\tskk Do not, my fair young ladies, said
he, stop your good ears against me.\tskk
Upon my word, honest man! said the
younger, we have no change.\tskk Then God
bless you, said the poor man, and multiply
those joys which you can give to others
without change!\tskk I observed the elder
sister put her hand into her pocket.\tskk
I’ll see, said she, if I have a sous.  A
sous! give twelve, said the supplicant;
Nature has been bountiful to you, be
bountiful to a poor man.

\tskk I would friend, with all my heart,
said the younger, if I had it.

My fair charitable! said he, addressing
himself to the elder\tskk what is it but
your goodness and humanity which makes
your bright eyes so sweet, that they
outshine the morning even in this dark
passage? and what was it which made the
Marquis de Santerre and his brother say so
much of you both as they just passed by?

The two ladies seemed much affected; and
impulsively, at the same time they both
put their hands into their pocket, and
each took out a twelve-sous piece.

The contest betwixt them and the poor
supplicant was no more\tskk it was
continued betwixt themselves, which of the
two should give the twelve-sous piece in
charity\tskk and, to end the dispute,
they both gave it together, and the man
went away.




\head{32pt}{48}{THE RIDDLE EXPLAINED}
\head{24pt}{48}{PARIS}

\vskip 6pt


\initial{I}{stepped} hastily after him: it
was the very man whose success in asking
charity of the women before the door of
the hotel had so puzzled me\tskk and I
found at once his secret, or at least the
basis of it\tskk ’twas flattery.

Delicious essence! how refreshing art thou
to Nature! how strongly are all its powers
and all its weaknesses on thy side! how
sweetly dost thou mix with the blood, and
help it through the most difficult and
tortuous passages to the heart!

The poor man, as he was not straiten’d for
time, had given it here in a larger dose:
’tis certain he had a way of bringing it
into a less form, for the many sudden
cases he had to do with in the streets:
but how he contrived to correct, sweeten,
concentre, and qualify it\tskk I vex not
my spirit with the enquiry\tskk it is
enough the beggar gained two twelve-sous
pieces\tskk and they can best tell the
rest, who have gained much greater matters
by it.




\head{32pt}{48}{PARIS}

\vskip 6pt


\dropcap{W}{e} get forwards in the world,
not so much by doing services, as
receiving them; you take a withering twig,
and put it in the ground; and then you
water it, because you have planted it.

Monsieur le Count de B\anon, merely
because he had done me one kindness in the
affair of my passport, would go on and do
me another, the few days he was at Paris,
in making me known to a few people of
rank; and they were to present me to
others, and so on.

I had got master of my \i{secret} just in
time to turn these honours to some little
account; otherwise, as is commonly the
case, I should have dined or supp’d a
single time or two round, and then, by
\i{translating} French looks and attitudes
into plain English, I should presently
have seen, that I had hold of the
\i{couvert} {652} of some more
entertaining guest; and in course should
have resigned all my places one after
another, merely upon the principle that I
could not keep them.\tskk As it was,
things did not go much amiss.

{652}  Plate, napkin, knife, fork and
spoon.

I had the honour of being introduced to
the old Marquis de B\tskk : in days of
yore he had signalized himself by some
small feats of chivalry in the \i{Cour
d’Amour}, and had dress’d himself out to
the idea of tilts and tournaments ever
since.\tskk The Marquis de B\tskk  wish’d
to have it thought the affair was
somewhere else than in his brain.  \lqq He
could like to take a trip to England,\rqq\ and
asked much of the English ladies.\tskk
Stay where you are, I beseech you,
Monsieur le Marquis, said I.\tskk \i{Les
Messieurs Anglois} can scarce get a kind
look from them as it is.\tskk The Marquis
invited me to supper.

Monsieur P\tskk , the farmer-general, was
just as inquisitive about our taxes.  They
were very considerable, he heard.\tskk If
we knew but how to collect them, said I,
making him a low bow.

I could never have been invited to Mons.
P\tskk ’s concerts upon any other terms.

I had been misrepresented to Madame de
Q\tskk  as an \i{esprit}.\tskk Madame de
Q\tskk  was an \i{esprit} herself: she
burnt with impatience to see me, and hear
me talk.  I had not taken my seat, before
I saw she did not care a sous whether I
had any wit or no\tskk I was let in, to
be convinced she had.  I call heaven to
witness I never once opened the door of my
lips.

Madame de V\tskk  vow’d to every creature
she met\tskk \lqq She had never had a more
improving conversation with a man in her
life.\rqq 

There are three epochas in the empire of a
French woman.\tskk She is coquette\tskk
then deist\tskk then \i{dévote}: the
empire during these is never lost\tskk
she only changes her subjects when
thirty-five years and more have unpeopled
her dominion of the slaves of love, she
re-peoples it with slaves of
infidelity\tskk and then with the slaves
of the church.

Madame de V\tskk  was vibrating betwixt
the first of those epochas: the colour of
the rose was fading fast away\tskk she
ought to have been a deist five years
before the time I had the honour to pay my
first visit.

She placed me upon the same sofa with her,
for the sake of disputing the point of
religion more closely.\tskk In short
Madame de V\tskk  told me she believed
nothing.\tskk I told Madame de V\tskk  it
might be her principle, but I was sure it
could not be her interest to level the
outworks, without which I could not
conceive how such a citadel as hers could
be defended\tskk that there was not a
more dangerous thing in the world than for
a beauty to be a deist\tskk that it was a
debt I owed my creed not to conceal it
from her\tskk that I had not been five
minutes sat upon the sofa beside her, but
I had begun to form designs\tskk and what
is it, but the sentiments of religion, and
the persuasion they had excited in her
breast, which could have check’d them as
they rose up?

We are not adamant, said I, taking hold of
her hand\tskk and there is need of all
restraints, till age in her own time
steals in and lays them on us.\tskk But my
dear lady, said I, kissing her hand\tskk
’tis too\tskk too soon.

I declare I had the credit all over Paris
of unperverting Madame de V\tskk .\tskk
She affirmed to Monsieur D\tskk  and the
Abbé M\tskk , that in one half hour I had
said more for revealed religion, than all
their Encyclopædia had said against
it.\tskk I was listed directly into Madame
de V\tskk ’s \i{coterie}\tskk and she put
off the epocha of deism for two years.

I remember it was in this \i{coterie}, in
the middle of a discourse, in which I was
shewing the necessity of a \i{first}
cause, when the young Count de Faineant
took me by the hand to the farthest corner
of the room, to tell me my \i{solitaire}
was pinn’d too straight about my
neck.\tskk It should be \i{plus badinant},
said the Count, looking down upon his
own\tskk but a word, Monsieur Yorick,
\i{to the wise}\tskk 

And \i{from the wise}, Monsieur le Count,
replied I, making him a bow\tskk \i{is
enough}.

The Count de Faineant embraced me with
more ardour than ever I was embraced by
mortal man.

For three weeks together I was of every
man’s opinion I met.\tskk \i{Pardi}! \i{ce
Monsieur Yorick a autant d’esprit que nous
autres}.\tskk \i{Il raisonne bien}, said
another.\tskk \i{C’est un bon enfant},
said a third.\tskk And at this price I
could have eaten and drank and been merry
all the days of my life at Paris; but
’twas a dishonest \i{reckoning}\tskk I
grew ashamed of it.\tskk It was the gain
of a slave\tskk every sentiment of honour
revolted against it\tskk the higher I
got, the more was I forced upon my
\i{beggarly system}\tskk the better the
\i{coterie}\tskk the more children of
Art\tskk I languish’d for those of
Nature: and one night, after a most vile
prostitution of myself to half a dozen
different people, I grew sick\tskk went
to bed\tskk order’d La Fleur to get me
horses in the morning to set out for
Italy.




\head{32pt}{48}{MARIA}
\head{24pt}{48}{MOULINES}

\vskip 6pt


\initial{I}{never} felt what the distress
of plenty was in any one shape till
now\tskk to travel it through the
Bourbonnois, the sweetest part of
France\tskk in the heyday of the vintage,
when Nature is pouring her abundance into
every one’s lap, and every eye is lifted
up\tskk a journey, through each step of
which Music beats time to \i{Labour}, and
all her children are rejoicing as they
carry in their clusters: to pass through
this with my affections flying out, and
kindling at every group before me\tskk
and every one of them was pregnant with
adventures.\tskk 

Just heaven!\tskk it would fill up twenty
volumes\tskk and alas! I have but a few
small pages left of this to crowd it
into\tskk and half of these must be taken
up with the poor Maria my friend, Mr.
Shandy, met with near Moulines.

The story he had told of that disordered
maid affected me not a little in the
reading; but when I got within the
neighbourhood where she lived, it returned
so strong into the mind, that I could not
resist an impulse which prompted me to go
half a league out of the road, to the
village where her parents dwelt, to
enquire after her.

’Tis going, I own, like the Knight of the
Woeful Countenance in quest of melancholy
adventures.  But I know not how it is, but
I am never so perfectly conscious of the
existence of a soul within me, as when I
am entangled in them.

The old mother came to the door; her looks
told me the story before she open’d her
mouth.\tskk She had lost her husband; he
had died, she said, of anguish, for the
loss of Maria’s senses, about a month
before.\tskk She had feared at first, she
added, that it would have plunder’d her
poor girl of what little understanding was
left\tskk but, on the contrary, it had
brought her more to herself\tskk still,
she could not rest.\tskk Her poor
daughter, she said, crying, was wandering
somewhere about the road.

Why does my pulse beat languid as I write
this? and what made La Fleur, whose heart
seem’d only to be tuned to joy, to pass
the back of his hand twice across his
eyes, as the woman stood and told it?  I
beckoned to the postilion to turn back
into the road.

When we had got within half a league of
Moulines, at a little opening in the road
leading to a thicket, I discovered poor
Maria sitting under a poplar.  She was
sitting with her elbow in her lap, and her
head leaning on one side within her
hand\tskk a small brook ran at the foot
of the tree.

I bid the postilion go on with the chaise
to Moulines\tskk and La Fleur to bespeak
my supper\tskk and that I would walk
after him.

She was dress’d in white, and much as my
friend described her, except that her hair
hung loose, which before was twisted
within a silk net.\tskk She had superadded
likewise to her jacket, a pale green
ribband, which fell across her shoulder to
the waist; at the end of which hung her
pipe.\tskk Her goat had been as faithless
as her lover; and she had got a little dog
in lieu of him, which she had kept tied by
a string to her girdle: as I looked at her
dog, she drew him towards her with the
string.\tskk \lqq Thou shalt not leave me,
Sylvio,\rqq\ said she.  I look’d in Maria’s
eyes and saw she was thinking more of her
father than of her lover, or her little
goat; for, as she utter’d them, the tears
trickled down her cheeks.

I sat down close by her; and Maria let me
wipe them away as they fell, with my
handkerchief.\tskk I then steep’d it in my
own\tskk and then in hers,\tskk and then
in mine\tskk and then I wip’d hers
again\tskk and as I did it, I felt such
undescribable emotions within me, as I am
sure could not be accounted for from any
combinations of matter and motion.

I am positive I have a soul; nor can all
the books with which materialists have
pester’d the world ever convince me to the
contrary.




\head{32pt}{48}{MARIA}

\vskip 6pt


\dropcap{W}{hen} Maria had come a little
to herself, I ask’d her if she remembered
a pale thin person of a man, who had sat
down betwixt her and her goat about two
years before?  She said she was unsettled
much at that time, but remembered it upon
two accounts\tskk that ill as she was,
she saw the person pitied her; and next,
that her goat had stolen his handkerchief,
and she had beat him for the theft\tskk
she had wash’d it, she said, in the brook,
and kept it ever since in her pocket to
restore it to him in case she should ever
see him again, which, she added, he had
half promised her.  As she told me this,
she took the handkerchief out of her
pocket to let me see it; she had folded it
up neatly in a couple of vine leaves, tied
round with a tendril\tskk on opening it,
I saw an S. marked in one of the corners.

She had since that, she told me, stray’d
as far as Rome, and walk’d round St.
Peter’s once\tskk and return’d back;\tskk
that she found her way alone across the
Apennines\tskk had travell’d over all
Lombardy, without money\tskk and through
the flinty roads of Savoy without
shoes\tskk how she had borne it, and how
she had got supported, she could not
tell\tskk but \i{God tempers the wind},
said Maria, \i{to the shorn lamb}.

Shorn indeed! and to the quick, said I:
and wast thou in my own land, where I have
a cottage, I would take thee to it, and
shelter thee: thou shouldst eat of my own
bread and drink of my own cup\tskk I
would be kind to thy Sylvio\tskk in all
thy weaknesses and wanderings I would seek
after thee and bring thee back\tskk when
the sun went down I would say my prayers:
and when I had done thou shouldst play thy
evening song upon thy pipe, nor would the
incense of my sacrifice be worse accepted
for entering heaven along with that of a
broken heart!

Nature melted within me, as I utter’d
this; and Maria observing, as I took out
my handkerchief, that it was steep’d too
much already to be of use, would needs go
wash it in the stream.\tskk And where will
you dry it, Maria? said I.\tskk I’ll dry
it in my bosom, said she\tskk ’twill do
me good.

And is your heart still so warm, Maria?
said I.

I touch’d upon the string on which hung
all her sorrows\tskk she look’d with
wistful disorder for some time in my face;
and then, without saying any thing, took
her pipe and play’d her service to the
Virgin.\tskk The string I had touched
ceased to vibrate\tskk in a moment or two
Maria returned to herself\tskk let her
pipe fall\tskk and rose up.

And where are you going, Maria? said
I.\tskk She said, to Moulines.\tskk Let us
go, said I, together.\tskk Maria put her
arm within mine, and lengthening the
string, to let the dog follow\tskk in
that order we enter’d Moulines.




\head{32pt}{48}{MARIA}
\head{24pt}{48}{MOULINES}

\vskip 6pt


\dropcap{T}{hough} I hate salutations and
greetings in the market-place, yet, when
we got into the middle of this, I stopp’d
to take my last look and last farewell of
Maria.

Maria, though not tall, was nevertheless
of the first order of fine forms\tskk
affliction had touched her looks with
something that was scarce earthly\tskk
still she was feminine\tskk and so much
was there about her of all that the heart
wishes, or the eye looks for in woman,
that could the traces be ever worn out of
her brain, and those of Eliza out of mine,
she should \i{not only eat of my bread and
drink of my own cup}, but Maria should lie
in my bosom, and be unto me as a daughter.

Adieu, poor luckless maiden!\tskk Imbibe
the oil and wine which the compassion of a
stranger, as he journeyeth on his way, now
pours into thy wounds\tskk the Being, who
has twice bruised thee, can only bind them
up for ever.




\head{32pt}{48}{THE BOURBONNNOIS}

\vskip 6pt


\dropcap{T}{here} was nothing from which I
had painted out for my self so joyous a
riot of the affections, as in this journey
in the vintage, through this part of
France; but pressing through this gate, of
sorrow to it, my sufferings have totally
unfitted me.  In every scene of festivity,
I saw Maria in the background of the
piece, sitting pensive under her poplar;
and I had got almost to Lyons before I was
able to cast a shade across her.

\tskk Dear Sensibility! source inexhausted
of all that’s precious in our joys, or
costly in our sorrows! thou chainest thy
martyr down upon his bed of straw\tskk and
’tis thou who lift’st him up to
Heaven!\tskk Eternal Fountain of our
feelings!\tskk ’tis here I trace thee\tskk
and this is thy \lqq \i{divinity which stirs
within me};\rqq \tskk not that, in some sad
and sickening moments, \lqq \i{my soul shrinks
back upon herself}, \i{and startles at
destruction};\rqq \tskk mere pomp of
words!\tskk but that I feel some generous
joys and generous cares beyond
myself\tskk all comes from thee,
great\tskk great SENSORIUM of the world!
which vibrates, if a hair of our heads but
falls upon the ground, in the remotest
desert of thy creation.\tskk Touch’d with
thee, Eugenius draws my curtain when I
languish\tskk hears my tale of symptoms,
and blames the weather for the disorder of
his nerves.  Thou giv’st a portion of it
sometimes to the roughest peasant who
traverses the bleakest mountains\tskk he
finds the lacerated lamb of another’s
flock.\tskk This moment I behold him
leaning with his head against his crook,
with piteous inclination looking down upon
it!\tskk Oh! had I come one moment sooner!
it bleeds to death!\tskk his gentle heart
bleeds with it.\tskk 

Peace to thee, generous swain!\tskk I see
thou walkest off with anguish\tskk but
thy joys shall balance it\tskk for, happy
is thy cottage\tskk and happy is the
sharer of it\tskk and happy are the lambs
which sport about you!




\head{32pt}{48}{THE SUPPER}

\vskip 6pt


\initial{A}{shoe} coming loose from the
fore foot of the thill-horse, at the
beginning of the ascent of mount Taurira,
the postilion dismounted, twisted the shoe
off, and put it in his pocket; as the
ascent was of five or six miles, and that
horse our main dependence, I made a point
of having the shoe fastened on again, as
well as we could; but the postilion had
thrown away the nails, and the hammer in
the chaise box being of no great use
without them, I submitted to go on.

He had not mounted half a mile higher,
when, coming to a flinty piece of road,
the poor devil lost a second shoe, and
from off his other fore foot.  I then got
out of the chaise in good earnest; and
seeing a house about a quarter of a mile
to the left hand, with a great deal to do
I prevailed upon the postilion to turn up
to it.  The look of the house, and of
every thing about it, as we drew nearer,
soon reconciled me to the disaster.\tskk
It was a little farm-house, surrounded
with about twenty acres of vineyard, about
as much corn\tskk and close to the house,
on one side, was a \i{potagerie} of an
acre and a half, full of everything which
could make plenty in a French peasant’s
house\tskk and, on the other side, was a
little wood, which furnished wherewithal
to dress it.  It was about eight in the
evening when I got to the house\tskk so I
left the postilion to manage his point as
he could\tskk and, for mine, I walked
directly into the house.

The family consisted of an old grey-headed
man and his wife, with five or six sons
and sons-in-law, and their several wives,
and a joyous genealogy out of them.

They were all sitting down together to
their lentil-soup; a large wheaten loaf
was in the middle of the table; and a
flagon of wine at each end of it promised
joy through the stages of the repast\tskk
’twas a feast of love.

The old man rose up to meet me, and with a
respectful cordiality would have me sit
down at the table; my heart was set down
the moment I enter’d the room; so I sat
down at once like a son of the family; and
to invest myself in the character as
speedily as I could, I instantly borrowed
the old man’s knife, and taking up the
loaf, cut myself a hearty luncheon; and,
as I did it, I saw a testimony in every
eye, not only of an honest welcome, but of
a welcome mix’d with thanks that I had not
seem’d to doubt it.

Was it this? or tell me, Nature, what else
it was that made this morsel so
sweet\tskk and to what magic I owe it,
that the draught I took of their flagon
was so delicious with it, that they remain
upon my palate to this hour?

If the supper was to my taste\tskk the
grace which followed it was much more so.




\head{32pt}{48}{THE GRACE}

\vskip 6pt


\dropcap{W}{hen} supper was over, the old
man gave a knock upon the table with the
haft of his knife, to bid them prepare for
the dance: the moment the signal was
given, the women and girls ran altogether
into a back apartment to tie up their
hair\tskk and the young men to the door
to wash their faces, and change their
sabots; and in three minutes every soul
was ready upon a little esplanade before
the house to begin.\tskk The old man and
his wife came out last, and placing me
betwixt them, sat down upon a sofa of turf
by the door.

The old man had some fifty years ago been
no mean performer upon the
\i{vielle}\tskk and at the age he was
then of, touch’d it well enough for the
purpose.  His wife sung now and then a
little to the tune\tskk then
intermitted\tskk and join’d her old man
again, as their children and
grand-children danced before them.

It was not till the middle of the second
dance, when, from some pauses in the
movements, wherein they all seemed to look
up, I fancied I could distinguish an
elevation of spirit different from that
which is the cause or the effect of simple
jollity.  In a word, I thought I beheld
\i{Religion} mixing in the dance\tskk
but, as I had never seen her so engaged, I
should have look’d upon it now as one of
the illusions of an imagination which is
eternally misleading me, had not the old
man, as soon as the dance ended, said,
that this was their constant way; and that
all his life long he had made it a rule,
after supper was over, to call out his
family to dance and rejoice; believing, he
said, that a cheerful and contented mind
was the best sort of thanks to heaven that
an illiterate peasant could pay\tskk 

Or a learned prelate either, said I.




\head{32pt}{48}{THE CASE OF DELICACY}

\vskip 6pt


\dropcap{W}{hen} you have gained the top
of Mount Taurira, you run presently down
to Lyons\tskk adieu, then, to all rapid
movements!  ’Tis a journey of caution; and
it fares better with sentiments, not to be
in a hurry with them; so I contracted with
a \i{voiturin} to take his time with a
couple of mules, and convoy me in my own
chaise safe to Turin, through Savoy.

Poor, patient, quiet, honest people! fear
not: your poverty, the treasury of your
simple virtues, will not be envied you by
the world, nor will your valleys be
invaded by it.\tskk Nature! in the midst
of thy disorders, thou art still friendly
to the scantiness thou hast created: with
all thy great works about thee, little
hast thou left to give, either to the
scythe or to the sickle\tskk but to that
little thou grantest safety and
protection; and sweet are the dwellings
which stand so shelter’d.

Let the way-worn traveller vent his
complaints upon the sudden turns and
dangers of your roads\tskk your
rocks\tskk your precipices;\tskk the
difficulties of getting up\tskk the
horrors of getting down\tskk mountains
impracticable\tskk and cataracts, which
roll down great stones from their summits,
and block his road up.\tskk The peasants
had been all day at work in removing a
fragment of this kind between St. Michael
and Madane; and, by the time my
\i{voiturin} got to the place, it wanted
full two hours of completing before a
passage could any how be gain’d: there was
nothing but to wait with patience\tskk
’twas a wet and tempestuous night; so that
by the delay, and that together, the
\i{voiturin} found himself obliged to put
up five miles short of his stage at a
little decent kind of an inn by the
roadside.

I forthwith took possession of my
bedchamber\tskk got a good fire\tskk
order’d supper; and was thanking heaven it
was no worse, when a \i{voiturin} arrived
with a lady in it and her servant maid.

As there was no other bed-chamber in the
house, the hostess\tskk without much
nicety, led them into mine, telling them,
as she usher’d them in, that there was
nobody in it but an English
gentleman\tskk that there were two good
beds in it, and a closet within the room
which held another.  The accent in which
she spoke of this third bed, did not say
much for it\tskk however, she said there
were three beds and but three people, and
she durst say, the gentleman would do
anything to accommodate matters.\tskk I
left not the lady a moment to make a
conjecture about it\tskk so instantly made
a declaration that I would do anything in
my power.

As this did not amount to an absolute
surrender of my bed-chamber, I still felt
myself so much the proprietor, as to have
a right to do the honours of it\tskk so I
desired the lady to sit down\tskk pressed
her into the warmest seat\tskk called for
more wood\tskk desired the hostess to
enlarge the plan of the supper, and to
favour us with the very best wine.

The lady had scarce warm’d herself five
minutes at the fire, before she began to
turn her head back, and give a look at the
beds; and the oftener she cast her eyes
that way, the more they return’d
perplexd\tskk I felt for her\tskk and for
myself: for in a few minutes, what by her
looks, and the case itself, I found myself
as much embarrassed as it was possible the
lady could be herself.

That the beds we were to lie in were in
one and the same room, was enough simply
by itself to have excited all this\tskk
but the position of them, for they stood
parallel, and so very close to each other
as only to allow space for a small wicker
chair betwixt them, rendered the affair
still more oppressive to us\tskk they
were fixed up moreover near the fire; and
the projection of the chimney on one side,
and a large beam which cross’d the room on
the other, formed a kind of recess for
them that was no way favourable to the
nicety of our sensations\tskk if anything
could have added to it, it was that the
two beds were both of them so very small,
as to cut us off from every idea of the
lady and the maid lying together; which in
either of them, could it have been
feasible, my lying beside them, though a
thing not to be wish’d, yet there was
nothing in it so terrible which the
imagination might not have pass’d over
without torment.

As for the little room within, it offer’d
little or no consolation to us: ’twas a
damp, cold closet, with a half dismantled
window-shutter, and with a window which
had neither glass nor oil paper in it to
keep out the tempest of the night.  I did
not endeavour to stifle my cough when the
lady gave a peep into it; so it reduced
the case in course to this
alternative\tskk That the lady should
sacrifice her health to her feelings, and
take up with the closet herself, and
abandon the bed next mine to her
maid\tskk or that the girl should take
the closet, \&c., \&c.

The lady was a Piedmontese of about
thirty, with a glow of health in her
cheeks.  The maid was a Lyonoise of
twenty, and as brisk and lively a French
girl as ever moved.\tskk There were
difficulties every way\tskk and the
obstacle of the stone in the road, which
brought us into the distress, great as it
appeared whilst the peasants were removing
it, was but a pebble to what lay in our
ways now.\tskk I have only to add, that it
did not lessen the weight which hung upon
our spirits, that we were both too
delicate to communicate what we felt to
each other upon the occasion.

We sat down to supper; and had we not had
more generous wine to it than a little inn
in Savoy could have furnish’d, our tongues
had been tied up, till necessity herself
had set them at liberty\tskk but the lady
having a few bottles of Burgundy in her
voiture, sent down her \i{fille de
chambre} for a couple of them; so that by
the time supper was over, and we were left
alone, we felt ourselves inspired with a
strength of mind sufficient to talk, at
least, without reserve upon our situation.
We turn’d it every way, and debated and
considered it in all kinds of lights in
the course of a two hours’ negotiation; at
the end of which the articles were settled
finally betwixt us, and stipulated for in
form and manner of a treaty of peace\tskk
and I believe with as much religion and
good faith on both sides as in any treaty
which has yet had the honour of being
handed down to posterity.

They were as follow\tskk 

First, as the right of the bed-chamber is
in Monsieur\tskk and he thinking the bed
next to the fire to be the warmest, he
insists upon the concession on the lady’s
side of taking up with it.

Granted, on the part of Madame; with a
proviso, That as the curtains of that bed
are of a flimsy transparent cotton, and
appear likewise too scanty to draw close,
that the Fille de Chambre shall fasten
up the opening, either by corking pins, or
needle and thread, in such manner as shall
be deem’d a sufficient barrier on the side
of Monsieur.

2dly.  It is required on the part of
Madame, that Monsieur shall lie the whole
night through in his \i{robe de chambre}.

Rejected: inasmuch as Monsieur is not
worth a \i{robe de chambre}; he having
nothing in his portmanteau but six shirts
and a black silk pair of breeches.

The mentioning the silk pair of breeches
made an intire change of the article\tskk
for the breeches were accepted as an
equivalent for the \i{robe de chambre};
and so it was stipulated and agreed upon,
that I should lie in my black silk
breeches all night.

3dly.  It was insisted upon and stipulated
for by the lady, that after Monsieur was
got to bed, and the candle and fire
extinguished, that Monsieur should not
speak one single word the whole night.

Granted; provided Monsieur’s saying his
prayers might not be deemed an infraction
of the treaty.

There was but one point forgot in this
treaty, and that was the manner in which
the lady and myself should be obliged to
undress and get to bed\tskk there was but
one way of doing it, and that I leave to
the reader to devise; protesting as I do
it, that if it is not the most delicate in
nature, ’tis the fault of his own
imagination\tskk against which this is
not my first complaint.

Now, when we were got to bed, whether it
was the novelty of the situation, or what
it was, I know not; but so it was, I could
not shut my eyes; I tried this side, and
that, and turn’d and turn’d again, till a
full hour after midnight; when Nature and
patience both wearing out\tskk O, my God!
said I.

\tskk You have broke the treaty, Monsieur,
said the lady, who had no more slept than
myself.\tskk I begg’d a thousand
pardons\tskk but insisted it was no more
than an ejaculation.  She maintained ’twas
an intire infraction of the treaty\tskk I
maintain’d it was provided for in the
clause of the third article.

The lady would by no means give up her
point, though she weaken’d her barrier by
it; for in the warmth of the dispute, I
could hear two or three corking pins fall
out of the curtain to the ground.

Upon my word and honour, Madame, said
I\tskk stretching my arm out of bed by
way of asseveration.\tskk 

(I was going to have added, that I would
not have trespassed against the remotest
idea of decorum for the world)\tskk 

But the Fille de Chambre hearing there
were words between us, and fearing that
hostilities would ensue in course, had
crept silently out of her closet, and it
being totally dark, had stolen so close to
our beds, that she had got herself into
the narrow passage which separated them,
and had advanced so far up as to be in a
line betwixt her mistress and me\tskk 

So that when I stretch’d out my
hand, I caught hold of the Fille de
Chambre’s

\stick{\tstrut 24pt\hfill\lsss{END}\enspace \lsss{OF}\enspace \lsss{VOL}.\enspace II.\hfill}

\end{document}
